\documentclass[11pt,a4paper]{article}

% Encoding and fonts
\usepackage[T1]{fontenc}
\usepackage[utf8]{inputenc}

% Mathematics
\usepackage{amsmath,amssymb}

% Page layout
\usepackage[margin=2cm]{geometry}

% Figures and captions
\usepackage{graphicx}
\usepackage{float}
\usepackage{caption}
\captionsetup{font=small,labelfont=bf}
% Shared publication visual system for generated vector-native result figures.
% Journal-quality TikZ/PGFPlots visual system.
% Add % Journal-quality TikZ/PGFPlots visual system.
% Add % Journal-quality TikZ/PGFPlots visual system.
% Add \input{figures/tikz/tikz_preamble.tex} once in the manuscript preamble.
\usepackage{tikz}
\usepackage{pgfplots}
\usepackage{pgfplotstable}
\usetikzlibrary{arrows.meta,backgrounds,calc,fit,matrix,patterns.meta,plotmarks,positioning}
\usepgfplotslibrary{colorbrewer,dateplot,fillbetween,groupplots,statistics}
\pgfplotsset{compat=1.18}

\definecolor{pubNavy}{HTML}{17365D}
\definecolor{pubBlue}{HTML}{2878B5}
\definecolor{pubTeal}{HTML}{2A9D8F}
\definecolor{pubGold}{HTML}{D79A22}
\definecolor{pubRose}{HTML}{C94C5C}
\definecolor{pubPurple}{HTML}{7656A7}
\definecolor{pubSlate}{HTML}{5F6B7A}
\definecolor{pubGray}{HTML}{9AA3AD}
\definecolor{pubLight}{HTML}{EDF1F5}
\definecolor{pubInk}{HTML}{18212B}

\pgfplotscreateplotcyclelist{publication cycle}{%
  pubNavy, thick\\
  pubTeal, thick, dashed\\
  pubGold, thick, dashdotted\\
  pubRose, thick, densely dashed\\
  pubPurple, thick, dotted\\
  pubSlate, thick, dash pattern=on 4pt off 1.5pt\\
}

\pgfplotsset{
  publication axis/.style={
    width=0.92\linewidth,
    height=0.42\linewidth,
    axis line style={draw=pubSlate!70, line width=0.35pt},
    tick style={draw=pubSlate!70, line width=0.35pt},
    tick label style={font=\scriptsize, text=pubInk},
    label style={font=\footnotesize, text=pubInk},
    title style={font=\footnotesize\bfseries, text=pubInk, yshift=0.6mm},
    legend style={font=\scriptsize, draw=none, fill=none,
      cells={anchor=west}, row sep=1pt},
    grid=major,
    grid style={draw=pubGray!24, line width=0.3pt},
    minor grid style={draw=pubGray!12, line width=0.2pt},
    clip mode=individual,
    cycle list name=publication cycle,
    every axis plot/.append style={line join=round, line cap=round},
  },
  publication group/.style={
    group style={horizontal sep=0.9cm, vertical sep=0.7cm},
  },
  publication point/.style={
    only marks, mark=*, mark size=2.2pt,
    draw=white, line width=0.45pt,
  },
  publication interval/.style={
    only marks, mark=*, mark size=2.3pt, thick,
    error bars/.cd, x dir=both, x explicit,
  },
}

\tikzset{
  publication note/.style={
    fill=white, fill opacity=0.94, text opacity=1,
    draw=pubGray!55, rounded corners=1.2pt,
    inner xsep=4pt, inner ysep=3pt,
    font=\scriptsize, text=pubInk, align=left,
  },
  publication warning/.style={
    publication note, draw=pubGold!80!black,
    fill=pubGold!8,
  },
}
 once in the manuscript preamble.
\usepackage{tikz}
\usepackage{pgfplots}
\usepackage{pgfplotstable}
\usetikzlibrary{arrows.meta,backgrounds,calc,fit,matrix,patterns.meta,plotmarks,positioning}
\usepgfplotslibrary{colorbrewer,dateplot,fillbetween,groupplots,statistics}
\pgfplotsset{compat=1.18}

\definecolor{pubNavy}{HTML}{17365D}
\definecolor{pubBlue}{HTML}{2878B5}
\definecolor{pubTeal}{HTML}{2A9D8F}
\definecolor{pubGold}{HTML}{D79A22}
\definecolor{pubRose}{HTML}{C94C5C}
\definecolor{pubPurple}{HTML}{7656A7}
\definecolor{pubSlate}{HTML}{5F6B7A}
\definecolor{pubGray}{HTML}{9AA3AD}
\definecolor{pubLight}{HTML}{EDF1F5}
\definecolor{pubInk}{HTML}{18212B}

\pgfplotscreateplotcyclelist{publication cycle}{%
  pubNavy, thick\\
  pubTeal, thick, dashed\\
  pubGold, thick, dashdotted\\
  pubRose, thick, densely dashed\\
  pubPurple, thick, dotted\\
  pubSlate, thick, dash pattern=on 4pt off 1.5pt\\
}

\pgfplotsset{
  publication axis/.style={
    width=0.92\linewidth,
    height=0.42\linewidth,
    axis line style={draw=pubSlate!70, line width=0.35pt},
    tick style={draw=pubSlate!70, line width=0.35pt},
    tick label style={font=\scriptsize, text=pubInk},
    label style={font=\footnotesize, text=pubInk},
    title style={font=\footnotesize\bfseries, text=pubInk, yshift=0.6mm},
    legend style={font=\scriptsize, draw=none, fill=none,
      cells={anchor=west}, row sep=1pt},
    grid=major,
    grid style={draw=pubGray!24, line width=0.3pt},
    minor grid style={draw=pubGray!12, line width=0.2pt},
    clip mode=individual,
    cycle list name=publication cycle,
    every axis plot/.append style={line join=round, line cap=round},
  },
  publication group/.style={
    group style={horizontal sep=0.9cm, vertical sep=0.7cm},
  },
  publication point/.style={
    only marks, mark=*, mark size=2.2pt,
    draw=white, line width=0.45pt,
  },
  publication interval/.style={
    only marks, mark=*, mark size=2.3pt, thick,
    error bars/.cd, x dir=both, x explicit,
  },
}

\tikzset{
  publication note/.style={
    fill=white, fill opacity=0.94, text opacity=1,
    draw=pubGray!55, rounded corners=1.2pt,
    inner xsep=4pt, inner ysep=3pt,
    font=\scriptsize, text=pubInk, align=left,
  },
  publication warning/.style={
    publication note, draw=pubGold!80!black,
    fill=pubGold!8,
  },
}
 once in the manuscript preamble.
\usepackage{tikz}
\usepackage{pgfplots}
\usepackage{pgfplotstable}
\usetikzlibrary{arrows.meta,backgrounds,calc,fit,matrix,patterns.meta,plotmarks,positioning}
\usepgfplotslibrary{colorbrewer,dateplot,fillbetween,groupplots,statistics}
\pgfplotsset{compat=1.18}

\definecolor{pubNavy}{HTML}{17365D}
\definecolor{pubBlue}{HTML}{2878B5}
\definecolor{pubTeal}{HTML}{2A9D8F}
\definecolor{pubGold}{HTML}{D79A22}
\definecolor{pubRose}{HTML}{C94C5C}
\definecolor{pubPurple}{HTML}{7656A7}
\definecolor{pubSlate}{HTML}{5F6B7A}
\definecolor{pubGray}{HTML}{9AA3AD}
\definecolor{pubLight}{HTML}{EDF1F5}
\definecolor{pubInk}{HTML}{18212B}

\pgfplotscreateplotcyclelist{publication cycle}{%
  pubNavy, thick\\
  pubTeal, thick, dashed\\
  pubGold, thick, dashdotted\\
  pubRose, thick, densely dashed\\
  pubPurple, thick, dotted\\
  pubSlate, thick, dash pattern=on 4pt off 1.5pt\\
}

\pgfplotsset{
  publication axis/.style={
    width=0.92\linewidth,
    height=0.42\linewidth,
    axis line style={draw=pubSlate!70, line width=0.35pt},
    tick style={draw=pubSlate!70, line width=0.35pt},
    tick label style={font=\scriptsize, text=pubInk},
    label style={font=\footnotesize, text=pubInk},
    title style={font=\footnotesize\bfseries, text=pubInk, yshift=0.6mm},
    legend style={font=\scriptsize, draw=none, fill=none,
      cells={anchor=west}, row sep=1pt},
    grid=major,
    grid style={draw=pubGray!24, line width=0.3pt},
    minor grid style={draw=pubGray!12, line width=0.2pt},
    clip mode=individual,
    cycle list name=publication cycle,
    every axis plot/.append style={line join=round, line cap=round},
  },
  publication group/.style={
    group style={horizontal sep=0.9cm, vertical sep=0.7cm},
  },
  publication point/.style={
    only marks, mark=*, mark size=2.2pt,
    draw=white, line width=0.45pt,
  },
  publication interval/.style={
    only marks, mark=*, mark size=2.3pt, thick,
    error bars/.cd, x dir=both, x explicit,
  },
}

\tikzset{
  publication note/.style={
    fill=white, fill opacity=0.94, text opacity=1,
    draw=pubGray!55, rounded corners=1.2pt,
    inner xsep=4pt, inner ysep=3pt,
    font=\scriptsize, text=pubInk, align=left,
  },
  publication warning/.style={
    publication note, draw=pubGold!80!black,
    fill=pubGold!8,
  },
}


% Tables
\usepackage{booktabs}

% Pseudocode
\usepackage{algorithm}
\usepackage{algorithmic}

% Bibliography
\usepackage[numbers,square,sort&compress]{natbib}

% Hyperlinks and cross-references
\usepackage[hidelinks]{hyperref}
\renewcommand{\eqref}[1]{Equation~\textup{(\ref{#1})}}

% Title information
\title{Neural Dynamic Vine Copula Pre-Training for Risk-Sensitive Reinforcement Learning in Dynamic Portfolio Optimization}
\author{Gabriël M.J.\ Gelissen and Fengyi Yuan\\
\small School of Science and Engineering (SSE), The Chinese University of Hong Kong, Shenzhen}
\date{}

\begin{document}

\maketitle

\begin{abstract}
Deep reinforcement learning provides a natural framework for sequential portfolio choice, but financial samples are short and joint tail dependence is difficult to learn from returns alone. We propose a synthetic-data-assisted portfolio controller that combines a fully dynamic, all-tree Student-$t$ D-vine with an LSTM-based Twin Delayed Deep Deterministic Policy Gradient (TD3) agent. Edge-specific neural networks update the D-vine correlations from lagged marginal and dependence information, while fixed edge degrees of freedom preserve identifiability. The fitted training-prefix model generates 1,000 synthetic 24-month episodes for pre-training; the agent is then fine-tuned once on 61 realised historical trajectories. The policy maximises multi-period terminal-wealth CRRA utility subject to a scenario-CVaR penalty, transaction and financing costs, and self-financing long--short constraints. In a frozen evaluation of seven US and Chinese equity indices and gold, the preregistered mean-weight ensemble of 20 independently trained policies is compared on the same realised path with equal weight, constrained shrinkage mean--variance, DCC--GARCH, static and rolling vines, and a dynamic neural vine optimiser without RL. Over the 22 periods satisfying the preregistered completeness rule, the ensemble achieves the highest observed Sharpe ratio (2.018) and lowest volatility (13.19\%), but ranks third in total return and annualised CRRA certainty-equivalent return. Its utility advantage over equal weight is positive but not statistically significant under paired circular-block bootstrap inference. Post-holdout explanatory ablations show that a coarse no-vine TD3 ensemble performs better descriptively and that historical fine-tuning adds little on this sample; these findings motivate a preregistered matched-seed causal decomposition of vine features, scenario-CVaR information, synthetic pre-training, recurrence, and reward shaping. The evidence supports competitive risk and implementation efficiency, rather than statistically proven superiority or an established incremental vine effect.
\end{abstract}

\medskip
\noindent\textit{Keywords:} Portfolio Optimisation, Dynamic Vine Copula, Deep Reinforcement Learning, TD3, LSTM, Synthetic Data, Tail Risk.

% Main sections
\section{Introduction}
\label{sec:introduction}

Portfolio optimization is a fundamental problem in financial decision-making, wherein investors seek to allocate capital among multiple assets to achieve an optimal balance between expected return and investment risk \cite{markowitz_portfolio_1952}. Since the seminal contribution of Markowitz, traditional portfolio theory has been built on the mean-variance framework, which models risk via the variance of portfolio returns. Despite its mathematical elegance and interpretability, the mean-variance framework rests on strong assumptions that are frequently violated in practice: normally distributed returns, constant market conditions, and fixed investment horizons \cite{campbell_econometrics_1998,demiguel_optimal_2009}. Financial markets are characterised by time-varying dynamics, volatility clustering, regime changes, and heavy-tailed return distributions---features that imply portfolio choices must be continuously recomputed as market conditions evolve \cite{campbell_econometrics_1998}. This necessitates sequential decision-making methods that can handle dynamically changing environments and manage uncertainty in a rigorous manner.

The modelling of multivariate dependencies in financial asset returns presents a fundamental challenge in portfolio optimization. Standard approaches, such as the Gaussian copula or the Dynamic Conditional Correlation (DCC) model of Engle \cite{engle_dynamic_2002}, impose restrictive elliptical dependence structures that may not adequately represent joint tail behaviour \cite{jondeau_copula-garch_2006,rodriguez_measuring_2007}. Vine copulas, introduced by Joe \cite{joe_multivariate_1997} and developed by Bedford and Cooke \cite{bedford_probability_2001,bedford_vinesnew_2002}, overcome this limitation by decomposing a high-dimensional joint distribution into a cascade of bivariate pair-copulas. The present implementation uses a fully dynamic Student-$t$ D-vine: all active tree edges are retained, degrees of freedom are fixed by edge for identifiability, and neural networks update edge correlations causally through time. This provides nonlinear, tail-dependent state information without repeatedly re-estimating a rolling vine or claiming unimplemented dynamic family switching.

So and Yeung \cite{so_vine-copula_2014} extended the vine copula framework to time-varying conditional dependence, demonstrating that dependence among financial assets varies significantly over time. Sahamkhadam and Stephan \cite{sahamkhadam_portfolio_2023} further showed that mixed vine copula models can improve risk reduction across stressed periods. Yet dependence estimation and portfolio control are distinct questions. Most vine-based allocation studies solve a fresh, myopic mean-variance, CVaR, or expected-utility problem at each date. Such optimisers do not directly learn a policy for the path-dependent interaction among current holdings, future rebalancing costs, leverage, and terminal wealth \cite{bernardi_portfolio_2018}.

Reinforcement learning has emerged as a powerful paradigm for sequential decision-making under uncertainty \cite{sutton_reinforcement_2018}. Casting portfolio allocation as a Markov decision process allows an agent to learn a state-contingent policy directly from repeated interactions with a market environment \cite{jang_deep_2023,jiang_deep_2017}. Benhamou et al. \cite{benhamou_bridging_2020} show how deep reinforcement learning can bridge Markowitz-style planning and sequential control. However, financial RL faces two practical difficulties: the available history is short relative to neural-network capacity, and an unconstrained return objective may reward unstable leverage, excessive turnover, or exposure to rare losses. These problems require both a defensible training-data protocol and an economic objective that accounts for downside risk and implementation.

Most portfolio RL states are built from returns, volatilities, and current holdings \cite{jiang_deep_2017,yu_model-based_2020}; the joint conditional distribution is typically implicit. Conversely, a copula generator can expose an agent to more paths, but synthetic data are useful only if their marginal, cross-sectional, tail, and temporal properties are validated. We therefore separate the roles of the dependence model. The neural D-vine (i) generates every episode used in synthetic pre-training, (ii) supplies explicit dynamic dependence features, and (iii) generates conditional scenarios from which the environment computes a policy-visible CVaR statistic and a reward penalty. Historical data ending before the locked final 24 decision periods are used for one fine-tuning pass; no synthetic ``fine-tuning'' sample and no evaluation return enter training.

This paper makes three methodological contributions. First, it integrates a fully dynamic, all-tree neural Student-$t$ D-vine with a recurrent TD3 controller, while keeping the dependence family and edge degrees of freedom fixed so that the claimed dynamics correspond exactly to the implemented correlation forecasts. Second, it formulates a self-financing long--short action map with net exposure one, gross exposure at most 1.5, individual weights in $[-0.20,0.60]$, proportional turnover costs, and explicit short- and cash-borrow charges. Dense reward increments telescope to terminal-wealth CRRA utility, so the objective is genuinely multi-period rather than a sequence of unrelated one-period utilities. Third, it imposes a leakage-resistant two-stage protocol: 1,000 synthetic 24-month episodes are used for pre-training, 61 realised training-prefix trajectories are used once for fine-tuning, and checkpoints must pass numerical, behavioural, and no-holdout gates before evaluation.

The empirical contribution is deliberately evidence-calibrated. The seven-asset adjusted-level panel runs from December 2013 to July 2026, and the final 24 monthly decisions are locked for evaluation. Twenty independently trained policies are combined by averaging target weights before returns and costs are recomputed. Against six causal benchmarks, the ensemble has the lowest observed volatility and highest observed Sharpe ratio over the 22 periods that satisfy the preregistered completeness rule, but it does not maximise terminal wealth or CRRA certainty-equivalent return and its utility effects are not statistically significant. Individual-seed dispersion is material, and much of the ensemble's risk advantage arises because opposing leveraged positions cancel in weight space. A frozen post-holdout explanatory ablation also finds no descriptive advantage over a coarse no-vine TD3 control and almost no increment from historical fine-tuning. We therefore report a mixed result, not a superiority claim.

These findings make causal attribution part of the research contribution. A prospective matched-seed analysis separately removes the direct neural-vine state, the scenario-CVaR observation, the CVaR reward term, synthetic pre-training, recurrence, and historical fine-tuning; it also compares the vine generator with a moving-block bootstrap and TD3 with DDPG, SAC, PPO, and A2C. All variants share the same realised return path, constraints, costs, environment-interaction budget, and ensemble construction. Because the original holdout has already been consumed, these mechanism results are explicitly post-holdout explanatory; fresh confirmation requires frozen non-overlapping walk-forward windows, an external investable panel, or future observations.

The remainder of this paper is organised as follows. Section \ref{sec:methodology} presents the neural D-vine, synthetic-data generator, LSTM-TD3 architecture, state and action design, multi-period objective, training protocol, benchmarks, and causal analysis contract. Section \ref{sec:data} describes the data and sample partition. Section \ref{sec:results} reports the frozen main evaluation, statistical inference, seed robustness, implementation effects, and explanatory ablations. Section \ref{sec:conclusion} concludes and identifies the evidence required for external and confirmatory validation.

%\section{Literature Review} 
\label{sec:lit_review}

% \lettrine[lraise=0.1, nindent=0.1em, slope=-.5em]{\textbf{P}}{ortfolio}


\subsection{Vine Copula Models for Portfolio Optimization}

\subsection{Deep Reinforcement Learning for Portfolio Management}

\subsection{Integration of Copula Methods and Machine Learning}

\subsection{Research Gap and Positioning}

\section{Methodology}
\label{sec:methodology}

This section presents the complete methodological framework underlying our vine-augmented deep reinforcement learning approach to dynamic portfolio optimisation. We begin by developing the D-vine copula model with neural network-driven time-varying parameters (\ref{subsec:3_Dvine_copula_model}), which serves both as the basis for the vine-augmented state representation and for generating synthetic training data. We then formalise the portfolio optimisation problem as a Markov decision process (\ref{subsec:3_markov_decision_process}), specifying the recurrent vine-augmented state, constrained continuous action space, and risk-sensitive multi-period reward. Next, we describe the LSTM-TD3 architecture (\ref{subsec:3_lstm_td3}) and the two-stage training strategy combining synthetic pre-training with historical fine-tuning (\ref{subsec:3_two_stage_training}). Finally, we define the benchmark models (\ref{subsec:3_benchmark_models}) and provide implementation details (\ref{subsec:3_implementation}).

\subsection{D-Vine Copula Model}
\label{subsec:3_Dvine_copula_model}

The modelling of multivariate dependencies in financial asset returns presents a fundamental challenge in portfolio optimization. Standard approaches, such as the Gaussian copula or the Dynamic Conditional Correlation (DCC) model, impose restrictive assumptions of linearity or elliptical dependence structures that fail to capture the asymmetric tail dependencies and non-linear relationships observed in financial markets during periods of stress. This section develops a flexible dependence framework based on vine copulas, which decompose a high-dimensional joint distribution into a cascade of bivariate pair-copulas, each of which can be selected from a broad family of parametric copula functions. The resulting framework accommodates complex, asymmetric, and tail-dependent structures while remaining computationally tractable for portfolios of moderate dimension.

\begin{figure}[H]
\centering
\fbox{\parbox{0.92\textwidth}{\centering
\textbf{Training prefix only}\par
Adjusted levels $\rightarrow$ daily log returns $\rightarrow$ diagnostically selected marginals
$\rightarrow$ frozen empirical PITs\par
$\Downarrow$\par
All-tree Student-$t$ D-vine $\rightarrow$ 21 edge-specific correlation networks
$\rightarrow$ causal monthly vine snapshots\par
$\Downarrow$\par
1,000 synthetic 24-month episodes $\rightarrow$ LSTM-TD3 pre-training
$\rightarrow$ one pass over 61 realised historical trajectories\par
$\Downarrow$\par
20 frozen seed policies $\rightarrow$ mean target weights $\rightarrow$ one common cost and constraint scorer
}}
\caption{Implemented Leakage-Controlled Neural-Vine and RL Pipeline}
\label{fig:vine_pipeline}
\end{figure}


\subsubsection{Copula Fundamentals and Sklar's Theorem}
\label{subsubsec:3_copula_fundamentals}

The theoretical foundation of all copula-based methods is Sklar's theorem \cite{sklar_fonctions_1959}, which establishes that any multivariate joint distribution can be expressed as a function of its univariate marginal distributions. Specifically, for an \( n \)-dimensional random vector \( \boldsymbol{Y} = (Y_1, \dots, Y_n) \) with joint distribution function \( F \) and marginal distribution functions \( F_1, \dots, F_n \), there exists an \( n \)-dimensional copula \( C \) such that

\begin{equation}
F(y_1, \dots, y_n) = C\{F_1(y_1), \dots, F_n(y_n)\},
\label{eq:sklar_joint}
\end{equation}

where the copula \( C: [0,1]^n \to [0,1] \) is itself a joint distribution function with uniform marginals. For continuous marginals, the joint density factorizes as

\begin{equation}
f(y_1, \dots, y_n) = c\{F_1(y_1), \dots, F_n(y_n)\} \cdot \prod_{i=1}^n f_i(y_i),
\label{eq:joint_density}
\end{equation}

where \( c \) is the copula density. The copula itself is defined as

\begin{equation}
C(u_1, \dots, u_n) = P(U_1 \leq u_1, \dots, U_n \leq u_n),
\label{eq:copula_def}
\end{equation}

with the transformed variables \( u_i = F_i(y_i) \sim \mathcal{U}[0,1] \) for \( i = 1, \dots, n \). This separation of marginal behaviour from dependence structure constitutes the central insight of the copula approach: the marginal dynamics of each asset can be modelled independently using univariate time-series models, while the joint dependence among assets is captured entirely by the copula function. 

The flexibility of the copula framework arises from the range of parametric pair-copula families available in principle. Table \ref{tab:copula_families} summarises the principal families used by the static and rolling benchmark estimators. The proposed dynamic generator deliberately restricts every active edge to the Student-$t$ family: the neural network changes dependence continuously through the correlation parameter, while the edge-specific degrees of freedom remain fixed. This restriction makes the dynamic claim identifiable and avoids discontinuous family switching.

\begin{table}[H]
\centering
\caption{D-Vine Copula Family Selection Criteria}
\label{tab:copula_families}
\begin{tabular}{l c c c}
\toprule
\textbf{Copula Family} & \textbf{Parameter Range} & \textbf{Tail Dependence} & \textbf{AIC/BIC Criterion} \\
\midrule
Gaussian & \(\rho \in (-1, 1)\) & No tail dependence & — \\
Student-t & \(\rho \in (-1, 1), \nu > 2\) & Symmetric tail dependence & — \\
Clayton & \(\theta \in (0, \infty)\) & Lower tail dependence & — \\
Gumbel & \(\theta \in [1, \infty)\) & Upper tail dependence & — \\
Frank & \(\theta \in \mathbb{R} \setminus \{0\}\) & No tail dependence & — \\
Joe & \(\theta \in [1, \infty)\) & Upper tail dependence & — \\
\bottomrule
\end{tabular}
\caption*{\footnotesize Notes: The table describes the available benchmark family set. The proposed dynamic D-vine uses Student-$t$ pair copulas on all active edges and does not switch families through time.}
\end{table}
\vspace{-15pt}

\subsubsection{Pair-Copula Construction (PCC)}
\label{subsubsec:3_pair_copula_construction}

For \( n > 2 \), specifying a valid \( n \)-dimensional copula directly is mathematically challenging, as not every combination of bivariate copulas yields a valid multivariate copula. The Pair-Copula Construction (PCC), introduced by Bedford and Cooke \cite{bedford_probability_2001,bedford_vinesnew_2002}, overcomes this limitation by decomposing the \( n \)-dimensional copula density into a product of \( n(n-1)/2 \) bivariate pair-copulas. This decomposition is hierarchical in nature: the first tree captures unconditional pairwise dependencies, the second tree captures dependencies conditional on one intermediate variable, and subsequent trees capture dependencies conditional on increasingly many variables.

For a D-vine with ordering \( 1-2-\dots-d \), the copula density is given by \eqref{eq:dvine_density}, where \( \ell \) indexes the tree level, \( k \) indexes the edge within each tree, and \( c_{k,k+\ell;k+1:k+\ell-1} \) denotes the conditional pair-copula density for variables \( k \) and \( k+\ell \) given the intermediate variables \( k+1, \dots, k+\ell-1 \). The functions \( \mathbb{C}_{k|k+1:k+\ell-1} \) represent the conditional distribution functions of \( U_k \) given the intermediate variables. This hierarchical structure allows the D-vine to capture higher-order dependencies that simple multivariate copulas, such as the Gaussian or Student-t, cannot represent.

\begin{equation}
\begin{aligned}
c_{1:d}(u_1, \dots, u_d) = \prod_{\ell=1}^{d-1} \prod_{k=1}^{d-\ell} c_{k,k+\ell;k+1:k+\ell-1}\{ &\mathbb{C}_{k|k+1:k+\ell-1}(u_k | \boldsymbol{u}_{k+1:k+\ell-1}), \\ &\mathbb{C}_{k+\ell|k+1:k+\ell-1}(u_{k+\ell} | \boldsymbol{u}_{k+1:k+\ell-1}) \}.
\label{eq:dvine_density}
\end{aligned}
\end{equation}

The conditional distributions in \eqref{eq:dvine_density} are evaluated using the so-called \( h \)-functions. For a pair-copula \( C_{ij;D} \), the \( h \)-function is defined as the partial derivative of the copula with respect to its second argument:

\begin{equation}
h_{i|j;D}(u_i | u_j) = \frac{\partial C_{ij;D}(u_i, u_j)}{\partial u_j}.
\label{eq:h_function}
\end{equation}

This function maps a pair of uniform variables to the conditional distribution of one given the other. The \( h \)-function enables the sequential evaluation of conditional distributions required for the vine decomposition:

\begin{equation}
\mathbb{C}_{i|D}(u_i | \boldsymbol{u}_D) = h_{i|j;D\setminus\{j\}}\{ \mathbb{C}_{i|D\setminus\{j\}}(u_i | \boldsymbol{u}_{D\setminus\{j\}}), \mathbb{C}_{j|D\setminus\{j\}}(u_j | \boldsymbol{u}_{D\setminus\{j\}}) \}.
\label{eq:conditional_dist}
\end{equation}

Similarly, for simulation from the vine copula, the inverse \( h \)-function is required:

\begin{equation}
u_i = h^{-1}_{i|j;D}(u_{i|D} | u_j).
\label{eq:h_inverse}
\end{equation}

The \( h \)-functions play a central role in both estimation and simulation. In estimation, they enable the sequential parameter estimation across vine trees: starting from Tree 1 (unconditional pair-copulas), parameters are estimated via maximum likelihood, the \( h \)-functions are computed, and the transformed variables are passed to Tree 2, where the process repeats. In simulation, the inverse \( h \)-functions are applied in reverse order to generate samples from the vine copula.

Figure \ref{fig:dvine_structure} illustrates the D-vine tree structure for a 7-asset portfolio, corresponding to the empirical application in this study. Tree 1 captures the unconditional pairwise dependencies among the seven assets, with the ordering \( 1-2-3-4-5-6-7 \). Tree 2 captures dependencies conditional on one intermediate variable, and this hierarchical structure continues through Tree 6, which captures the highest-order conditional dependence.

\begin{figure}[H]
\centering
\includegraphics[width=0.9\textwidth]{figures/d_vine_structure.pdf}
\caption{D-Vine Tree Structure for 7 Assets}
\label{fig:dvine_structure}
\end{figure}

\subsubsection{Marginal Models and Causal Probability Integral Transforms}
\label{subsubsec:3_marginal_models}

Before the copula can be estimated, each marginal return process must absorb serial dependence, volatility clustering, leverage effects, and heavy tails. Marginals are fitted once on the training prefix; their parameters and empirical residual distributions are then frozen. No evaluation observation can alter either the fitted model or the ranks used to construct pseudo-observations.

For each asset, the candidate conditional mean is either constant or AR(1):

\begin{equation}
r_{i,t}=\mu_i+\phi_i r_{i,t-1}+\varepsilon_{i,t},\qquad
\varepsilon_{i,t}=\sigma_{i,t}z_{i,t},
\label{eq:ar1_mean}
\end{equation}

where $\phi_i=0$ for the constant-mean candidate. The candidate volatility grid contains symmetric GARCH, GJR-GARCH, and exponential GARCH specifications, each with order $(1,1)$ and Student-$t$ or skewed Student-$t$ innovations. The symmetric candidate is

\begin{equation}
\sigma_{i,t}^2=\omega_i+\alpha_i\varepsilon_{i,t-1}^2+\beta_i\sigma_{i,t-1}^2,
\label{eq:garch11}
\end{equation}

with the usual positivity and covariance-stationarity conditions

\begin{equation}
\omega_i>0,\qquad \alpha_i\geq0,\qquad\beta_i\geq0,
\qquad\alpha_i+\beta_i<1.
\label{eq:garch_constraints}
\end{equation}

The GJR candidate adds an indicator-weighted negative-shock term and the eGARCH candidate evolves log variance. Every converged candidate is tested for remaining autocorrelation in both $z_{i,t}$ and $z_{i,t}^2$ at lags 5, 10, and 20. A candidate enters the BIC selection pool only if the minimum p-value in both diagnostic families is at least 0.01. If no one-component GARCH candidate passes, a causal two-component EWMA grid with short- and long-horizon variance filters is evaluated under the same residual gate. The run fails closed if no admissible marginal remains.

The standardised residuals are

\begin{equation}
z_{i,t}=\frac{\varepsilon_{i,t}}{\sigma_{i,t}},
\label{eq:standardized_residuals}
\end{equation}

and the probability integral transform required by the copula is implemented with training-residual mid-ranks:

\begin{equation}
u_{i,t}=\widehat F_{i,\mathrm{train}}(z_{i,t}).
\label{eq:pit}
\end{equation}

Values are clipped away from zero and one. For later dates, fixed-parameter filtering produces $z_{i,t}$ and the frozen training empirical distribution maps it to $u_{i,t}$; hence evaluation ranks cannot leak backward into model fitting. The likelihood for a candidate marginal has the generic form

\begin{equation}
\mathcal{L}_i(\boldsymbol\theta_i)=\sum_{t\in\mathcal T_{\mathrm{train}}}
\left[\log f_{\eta_i}(z_{i,t})-\log\sigma_{i,t}\right],
\label{eq:garch_likelihood}
\end{equation}

where $f_{\eta_i}$ is the selected innovation density and $\boldsymbol\theta_i$ includes the corresponding mean, variance, shape, and skew parameters. Hybrid numerical optimisation is used, non-converged fits are discarded, and BIC is compared only after the residual gate. This separation allows each asset to retain its own marginal dynamics while the D-vine models the cross-sectional dependence of the filtered uniforms.

\subsubsection{Dynamic D-Vine with Neural Network Parametrisation}
\label{subsubsec:3_dynamic_vine_with_NN_param}

The vine copula framework described thus far is static, whereas financial dependence evolves through time \cite{so_vine-copula_2014}. The implemented model is a fully dynamic Student-$t$ D-vine. A Kendall's-tau path selects a fixed D-vine order on the training prefix, all $d(d-1)/2=21$ edges are retained for $d=7$, and every edge is fitted with a Student-$t$ pair copula. Edge-specific degrees of freedom $\nu_e$ are fixed after training-prefix estimation; an edge-specific neural network forecasts the time-varying correlation $\rho_{e,t}$.

For edge $e$, the eight-dimensional causal feature vector is

\begin{equation}
\boldsymbol x_{e,t}=\left(u_{1,t-1},u_{2,t-1},z_{1,t-1},z_{2,t-1},
\log\sigma_{1,t-1},\log\sigma_{2,t-1},\rho_{e,t-1},s_{e,t-1}\right)^{\top},
\label{eq:tau_evolution}
\end{equation}

where $s_{e,t-1}$ is the lagged Student-$t$ copula score. For higher trees, the two uniform inputs are causal conditional pseudo-observations obtained recursively from the preceding trees. A two-hidden-layer network with widths 16 and 8 maps these features to a bounded correlation:

\begin{equation}
\rho_{e,t}=\tanh\!\left[\operatorname{atanh}(\bar\rho_e)+a_e
g_e(\boldsymbol x_{e,t};\boldsymbol w_e)\right],
\label{eq:nn_parameter}
\end{equation}

where $\bar\rho_e$ is the static training-prefix estimate and $a_e$ limits dynamic departures. The network is trained directly by the Student-$t$ copula negative log-likelihood, rather than by regression to rolling-window estimates:

\begin{equation}
\mathcal{L}_{e}(\boldsymbol w_e)=-\sum_{t\in\mathcal T_{\mathrm{fit}}}
\log c_{t,\nu_e}\!\left(u_{1,t},u_{2,t};\rho_{e,t}(\boldsymbol w_e)\right).
\label{eq:nn_objective}
\end{equation}

The final 20\% of the training sequence is kept in chronological order for early stopping, with patience 50. Networks are initialised at the static edge estimate, so a weak dynamic signal reduces to the fitted backbone rather than an arbitrary dependence path. The resulting model is estimated once on the training prefix and replayed causally; it is not a rolling-window vine.

Table \ref{tab:copula_families_definitions} records the broader family set relevant to conventional vine benchmarks. In the proposed neural model, only the Student-$t$ row is active. Consequently, lower and upper tail dependence are symmetric conditional on an edge, but remain time-varying through $\rho_{e,t}$.

\begin{table}[H]
\centering
\caption{Bivariate Copula Families Used in the Framework and Benchmarks}
\resizebox{\textwidth}{!}{%
\begin{tabular}{l c c c c}
\toprule
\textbf{Family} & \multicolumn{1}{c}{$\mathbf{C(u,v)}$} & \textbf{Parameter Range} & $\mathbf{\lambda_L}$ & $\mathbf{\lambda_U}$ \\
\midrule
Gaussian & $\Phi_\rho(\Phi^{-1}(u), \Phi^{-1}(v))$ & $\rho \in (-1,1)$ & 0 & 0 \\
Student-$t$ & $T_{\rho,\nu}(T_\nu^{-1}(u), T_\nu^{-1}(v))$ & $\rho \in (-1,1), \nu > 2$ & $2t_{\nu+1}\left(-\sqrt{\frac{(\nu+1)(1-\rho)}{1+\rho}}\right)$ & same as $\lambda_L$ \\
Clayton & $\left(u^{-\theta} + v^{-\theta} - 1\right)^{-1/\theta}$ & $\theta > 0$ & $2^{-1/\theta}$ & 0 \\
Gumbel & $\exp\left\{-\left[(-\ln u)^\theta + (-\ln v)^\theta\right]^{1/\theta}\right\}$ & $\theta \geq 1$ & 0 & $2 - 2^{1/\theta}$ \\
Frank & $-\frac{1}{\theta}\ln\left\{1 + \frac{(e^{-\theta u} - 1)(e^{-\theta v} - 1)}{e^{-\theta} - 1}\right\}$ & $\theta \neq 0$ & 0 & 0 \\
Joe & $1 - \left\{(1-u)^\theta + (1-v)^\theta - (1-u)^\theta(1-v)^\theta\right\}^{1/\theta}$ & $\theta \geq 1$ & 0 & $2 - 2^{1/\theta}$ \\
\bottomrule
\end{tabular}%
}
\label{tab:copula_families_definitions}
\caption*{\footnotesize Notes: $\Phi_\rho$ denotes the bivariate normal CDF with correlation $\rho$; $\Phi^{-1}$ is the inverse standard normal CDF. $T_{\rho,\nu}$ denotes the bivariate Student-$t$ CDF with correlation $\rho$ and degrees of freedom $\nu$; $T_\nu^{-1}$ is the inverse univariate Student-$t$ CDF. $\lambda_L$ and $\lambda_U$ denote lower and upper tail dependence, respectively. For the Student-$t$ copula, lower and upper tail dependence are symmetric.}
\end{table}
\vspace{-15pt}

The tail dependence coefficients quantify limiting joint extremes. In the implemented state, each active edge contributes its correlation, implied Kendall's $\tau$, and symmetric Student-$t$ tail coefficient. These 63 features inform the RL policy, and the same causal edge snapshots generate conditional scenarios. Dynamic family selection and asymmetric upper/lower coefficients are therefore not claimed for the principal model.


\subsubsection{Synthetic Data Generation}
\label{subsubsec:3_synthetic_data_generation}

A central contribution of this study is the use of the training-prefix neural D-vine as a controlled data generator. The generator is not assumed to reproduce every feature of the market merely because its marginals pass a diagnostic. It is accepted only after separate marginal, Pearson-correlation, lower-tail co-exceedance, and temporal-dependence checks compare simulated and historical training-prefix data. These diagnostics are treated as generator validation, not as evidence of investment performance.

The daily marginal filters first produce training-only pseudo-observations. These are aggregated to monthly holding-period uniforms, and a fitted univariate Student-$t$ Markov copula for each asset supplies serial and tail persistence. Conditional innovations retain the contemporaneous neural D-vine dependence. At month $t$, forward simulation draws

\begin{equation}
\boldsymbol v_t\sim C^{\mathrm{NN\mbox{-}vine}}_t,\qquad
\boldsymbol u_t=\mathcal S(\boldsymbol v_t,\boldsymbol u_{t-1}),
\label{eq:dvine_sample}
\end{equation}

where $\mathcal S$ denotes the seven fitted serial-copula transitions. The resulting uniforms are mapped through continuous empirical quantiles of training-prefix monthly log returns:

\begin{equation}
\widetilde r_{i,t}=\widehat Q_{i,\mathrm{train}}(u_{i,t}),
\label{eq:inverse_returns}
\end{equation}

using midpoint plotting positions to avoid excess mass beyond an isolated minimum. A single Fisher-$z$ dependence multiplier is calibrated by simulated method of moments against the 21 training-prefix Pearson correlations. This parsimonious correction preserves edge signs and dynamic ordering and avoids pair-specific calibration to 21 noisy targets.

The realised training bundle contains $M=1{,}000$ synthetic episodes. Every episode has a 30-month causal burn-in followed by $T'=24$ monthly decisions. For each decision, one draw supplies the realised synthetic transition and 512 additional draws supply the conditional loss distribution used for CVaR:

\begin{equation}
\mathcal D_{\mathrm{synth}}=
\left\{\left(S^{(m)}_{0},\{\widetilde{\boldsymbol R}^{(m)}_{t,j}\}_{j=0}^{512}\right)_{t=1}^{24}\right\}_{m=1}^{1000}.
\label{eq:synth_paths}
\end{equation}

The starting dependence snapshot is sampled only from causal sequences that fit within the training prefix. No historical realised return is inserted into a synthetic episode. Conversely, fine-tuning episodes use the actual permissible historical return as row zero and use simulated rows only as ex-ante CVaR scenarios. Thus all synthetic episodes serve pre-training, all permissible realised history serves fine-tuning, and the locked final 24 decision periods serve evaluation only.

Algorithm \ref{alg:dvine} summarises the implemented generator. The D-vine and marginal models are fitted once; monthly snapshots are then replayed without rolling re-estimation.

\begin{algorithm}[H]
\caption{Neural D-Vine Estimation and Synthetic Episode Generation}
\label{alg:dvine}
\begin{algorithmic}[1]
\STATE \textbf{Input:} Training-prefix daily returns; $N=7$; $M=1{,}000$; $T'=24$; burn-in $L=30$
\STATE \textbf{Output:} Frozen marginal models, causal D-vine snapshots, synthetic episodes
\STATE Fit and diagnostically gate the marginal candidate grid on the training prefix
\STATE Construct frozen empirical PITs and monthly serial-copula innovations
\STATE Select one D-vine order and fit an all-tree Student-$t$ backbone
\FOR{each of the 21 edges}
    \STATE Train the correlation network by copula likelihood with chronological validation
\ENDFOR
\STATE Calibrate one training-only Fisher-$z$ dependence scale by simulated moments
\STATE Build the causal sequence $\{C_t^{\mathrm{NN\mbox{-}vine}\}}$
\FOR{$m=1,\ldots,M$}
    \STATE Sample a valid $L+T'$ snapshot sequence
    \FOR{$t=1,\ldots,L+T'$}
        \STATE Draw 513 D-vine innovations and apply the serial transitions
        \STATE Map uniforms through training-prefix monthly empirical quantiles
    \ENDFOR
    \STATE Store burn-in context and the 24 realised/scenario matrices
\ENDFOR
\STATE Validate marginal, correlation, tail, and temporal fidelity; write immutable hashes
\RETURN the frozen training bundle
\end{algorithmic}
\end{algorithm}







\subsection{Markov Decision Process (MDP) Formulation}
\label{subsec:3_markov_decision_process}

To model portfolio optimisation as a sequential decision-making problem, we must formalise the investment process as a Markov Decision Process (MDP) \cite{sutton_reinforcement_2018}. The MDP is defined by the tuple $(\mathcal{S}, \mathcal{A}, \mathcal{P}, \mathcal{R}, \gamma)$, where $\mathcal{S}$ is the state space, $\mathcal{A}$ is the action space, $\mathcal{P}$ is the state transition kernel, $\mathcal{R}$ is the reward function, and $\gamma \in (0,1]$ is the discount factor. At each rebalancing period $t$, the agent observes the current market state $s_t \in \mathcal{S}$, selects a portfolio allocation $a_t \in \mathcal{A}$, receives a reward $r_t = \mathcal{R}(s_t, a_t)$, and transitions to a new state $s_{t+1}$ according to the market dynamics $\mathcal{P}(s_{t+1} | s_t, a_t)$. The objective is to learn a policy $\pi: \mathcal{S} \to \mathcal{A}$ that maximises the expected cumulative discounted reward $\mathbb{E}[\sum_{t=0}^{T} \gamma^t r_t]$. Table \ref{tab:mdp_components} summarises the MDP components, which are detailed in the following subsections.

\begin{table}[H]
\centering
\caption{Markov Decision Process Components for Portfolio Optimisation}
\label{tab:mdp_components}
\begin{tabular}{p{2.5cm} p{10cm}}
\toprule
\textbf{Component} & \textbf{Definition} \\
\midrule
State $S_t$ & A 30-observation sequence containing wealth, monthly returns, volatilities, scenario CVaR, prior holdings and exposures, and all D-vine edge features \\
Action $a_t$ & Self-financing long-short weights with net exposure 1, gross exposure at most 1.5, and asset limits of $+0.60/-0.20$ \\
Reward $r_t$ & Increment in terminal-wealth CRRA utility minus scenario-based CVaR; transaction and financing costs enter wealth and CVaR directly \\
Transition $P$ & Stochastic market dynamics governing asset returns (unknown to the agent, learned through experience) \\
Discount $\gamma$ & Discount factor controlling the trade-off between immediate and future rewards \\
\bottomrule
\end{tabular}
\end{table}

\subsubsection{State Space Design}
\label{subsubsec:3_state_space_design}

The state design determines the information available to every actor and critic update. The implementation distinguishes the current observation $o_t$ from the recurrent state $S_t$. Each observation contains normalised wealth, the most recently realised vector of monthly log returns, annualised rolling volatilities, the preceding action, its gross and net exposures, the most recently estimated scenario CVaR, and the complete D-vine feature vector. Rolling Pearson correlations are not supplied separately. The recurrent input is the ordered sequence
%
\begin{equation}
S_t=(o_{t-L+1},\ldots,o_t), \qquad L=30,
\label{eq:recurrent_state}
\end{equation}
%
so temporal information is represented by 30 monthly observations rather than by embedding 30 daily returns inside one vector.

For asset $i$, the return component is the monthly log gross return $r_{i,t}=\log G_{i,t}$, expressed in percentage points before entering the network. Volatility is the annualised root mean square of the latest 20 monthly log returns,
%
\begin{equation}
\sigma_{i,t}=100\sqrt{12\left(\frac{1}{20}\sum_{j=0}^{19}r_{i,t-j}^{2}\right)},
\label{eq:volatility}
\end{equation}
%
which supplies a causal measure of changing return scale. Every generated training episode also contains exactly 30 burn-in returns and corresponding vine states preceding its first decision. This prevents the LSTM history from being formed by repeating the reset observation and ensures that the first action is conditioned only on information available before that action.

The D-vine block contains three features for every unconditional and conditional pair-copula edge: Kendall's tau and lower- and upper-tail dependence. For an edge copula $C_{e,t}$,
%
\begin{equation}
\tau_{e,t}=4\int_0^1\int_0^1 C_{e,t}(u,v)\,dC_{e,t}(u,v)-1,
\qquad
(\lambda^L_{e,t},\lambda^U_{e,t}),
\label{eq:vine_state_features}
\end{equation}
%
where $\lambda^L_{e,t}=\lambda^U_{e,t}$ is evaluated from the edge's dynamic Student-$t$ correlation and fixed degrees of freedom. The coarse no-vine ablation preserves the same observation dimension but sets this vector, together with the policy-visible scenario-CVaR channel, to zero. The causal analysis additionally masks these two channels separately, so direct dependence features are not conflated with a risk statistic derived from vine scenarios.

Portfolio context consists of $V_t/V_0$, the drifted pre-trade weight vector $\widetilde w_{t^-}$, $\lVert\widetilde w_{t^-}\rVert_1$, and $\mathbf{1}^{\top}\widetilde w_{t^-}$. The reset allocation is the neutral equal-weight portfolio, not a zero-holdings or all-cash state. The scenario-CVaR component is the mean loss in the worst 5\% of the 512 joint monthly return scenarios used at the preceding step. Accordingly,
%
\begin{equation}
o_t=\left[\frac{V_t}{V_0},\ 100r_t,\ \sigma_t,\ 100\operatorname{CVaR}_{0.95,t},\ \widetilde w_{t^-},\ \lVert\widetilde w_{t^-}\rVert_1,\ \mathbf{1}^{\top}\widetilde w_{t^-},\ \theta^{\mathrm{vine}}_t\right]^{\top}.
\label{eq:state_vector}
\end{equation}
%
With $N=7$ assets and $3N(N-1)/2=63$ vine features, a single observation has
%
\begin{equation}
d_o=1+N+N+1+N+2+\frac{3N(N-1)}{2}=88
\label{eq:state_dimension}
\end{equation}
%
components, and each recurrent network receives a tensor of dimension $30\times88$ per decision.

\subsubsection{Action Space}
\label{subsubsec:3_action_space}

At each monthly rebalancing date, the actor produces $N+1$ raw outputs: one cross-sectional score per asset and one leverage logit. Scores are smoothly bounded as $d_{i,t}=b\tanh(z_{i,t}/b)$ with $b=5$, while the leverage gate is $g_t=\operatorname{sigmoid}(z_{N+1,t})\in(0,1)$. The feasible mandate is
%
\begin{equation}
\mathcal{A}=\left\{w\in\mathbb{R}^{N}:\ \mathbf{1}^{\top}w=1,\ \lVert w\rVert_1\leq1.5,\ -0.20\leq w_i\leq0.60\right\}.
\label{eq:action_space}
\end{equation}
%
Thus, the implemented policy is self-financing and long-short rather than long-only.

The two lowest-ranked asset scores form a disjoint short support and the remaining assets form the long support. For net exposure one and maximum gross exposure 1.5, the gate determines the short and long budgets as
%
\begin{equation}
B_t^- = 0.25g_t, \qquad B_t^+ = 1+B_t^-.
\label{eq:book_budgets}
\end{equation}
%
Within each support, an interior logistic capped-simplex map with temperature 1.5 converts scores to probability vectors $p_t^+$ and $p_t^-$. Its caps are $0.60/B_t^+$ for the long probabilities and $0.20/B_t^-$ for the short probabilities. Portfolio weights are then
%
\begin{equation}
w_t=B_t^+p_t^+-B_t^-p_t^-, \qquad
\lVert w_t\rVert_1=1+0.5g_t.
\label{eq:constrained_action_map}
\end{equation}
%
Because the supports do not overlap, the gate is identifiable with realised gross exposure, and the net, gross, and individual-position constraints hold by construction. The smooth interior map retains gradients near the position limits, unlike a clipped softmax or Euclidean projection that can create exact boundary plateaus.

\subsubsection{Reward Function}
\label{subsubsec:3_reward_function}

The reward implements a multi-period terminal-wealth objective with dense learning signals. Let $G_t$ denote the vector of realised holding-period gross asset returns. For newly selected target weights $w_t$, the pre-cost growth factor and economically implementable turnover are
%
\begin{equation}
g_{p,t}=1+w_t^{\top}(G_t-\mathbf{1}),
\qquad \operatorname{TO}_t=\lVert w_t-\widetilde w_{t^-}\rVert_1.
\label{eq:portfolio_return}
\end{equation}
%
Transaction and financing costs are applied multiplicatively rather than subtracted as a second reward term. With proportional cost $\kappa=0.001$, annual short-borrow rate $c_s=0.03$, annual cash-borrow rate $c_c=0.02$, and holding-period year fraction $\Delta_t$, the cost multiplier is
%
\begin{equation}
m_t=\exp\left[-\kappa\operatorname{TO}_t-\Delta_t\{c_s\sum_i(-w_{i,t})_+ + c_c(\mathbf{1}^{\top}w_t-1)_+\}\right],
\qquad V_{t+1}=V_t\max(g_{p,t}m_t,10^{-12}).
\label{eq:cost_adjusted_wealth}
\end{equation}
%
After the realised move, the next pre-trade risky holdings are $\widetilde w_{(t+1)^-}=w_t\odot G_t/g_{p,t}$. The same target weights and costs are applied to 512 joint scenarios generated from the contemporaneous vine. For scenario loss $\ell_t^{(m)}=1-g_{p,t}^{(m)}m_t$, $\operatorname{CVaR}_{0.95,t}$ is the mean of the largest $\lceil0.05\times512\rceil=26$ losses.

CRRA utility is defined on the wealth ratio, $U_{\gamma_U}(x)=(x^{1-\gamma_U}-1)/(1-\gamma_U)$ for $\gamma_U\neq1$, with $\gamma_U=2$. The implemented reward is
%
\begin{equation}
r_t=U_{\gamma_U}\!\left(\frac{V_{t+1}}{V_0}\right)-U_{\gamma_U}\!\left(\frac{V_t}{V_0}\right)-\lambda\operatorname{CVaR}_{0.95,t},
\qquad \lambda=0.05.
\label{eq:reward_function}
\end{equation}
%
The utility increments telescope across the 24-month episode to terminal-wealth utility, while still providing a reward after every decision. The agent discount is therefore fixed at one. Turnover, short-borrow costs, and any cash-borrow costs affect both realised wealth and the forward loss distribution; they are not counted again through a separate additive turnover penalty.









\subsection{LSTM-TD3 Architecture}
\label{subsec:3_lstm_td3}

The current agent is a recurrent implementation of Twin Delayed Deep Deterministic Policy Gradient (TD3) \cite{fujimoto_addressing_2018}, which extends DDPG \cite{lillicrap_continuous_2015} for continuous control. It retains a deterministic actor but replaces the single critic with two independently parameterised critics, uses the smaller target estimate, smooths target actions with clipped Gaussian noise, and updates the actor less frequently than the critics. These changes address value overestimation and rapidly moving policy targets. The recurrence is not a shared feature extractor: the actor and each critic own a separate LSTM encoder and therefore learn task-specific temporal representations.

Figure \ref{fig:lstm_td3_architecture} illustrates the implemented architecture. A 30-observation vine-augmented sequence is normalised feature-wise and processed by a one-layer LSTM. The actor converts its last temporal representation into $N$ asset scores and one leverage logit, which are passed through the constrained action map in Equation \eqref{eq:constrained_action_map}. Each critic processes the state sequence through its own recurrent encoder, concatenates the last hidden representation with the realised constrained action, and estimates a scalar action value.

\begin{figure}[H]
\centering
\input{figures/tikz/figure_m01_lstm_td3_architecture.tex}
\caption{Implemented LSTM-TD3 Architecture and Update Loop}
\label{fig:lstm_td3_architecture}
\caption*{\footnotesize Notes: The actor, critic 1, and critic 2 have separate LayerNorm--LSTM encoders. The actor's $N+1$ raw outputs are mapped to capped disjoint long and short books. TD3 trains both critics against the minimum smoothed target value and updates the actor and all target networks every second critic update.}
\end{figure}

\subsubsection{LSTM Temporal Encoding}
\label{subsubsec:3_lstm_temporal_encoding}

Financial returns and dependence states can exhibit persistence, regime shifts, and volatility clustering. Each recurrent network therefore processes the causal sequence $S_t$ with an LSTM \cite{hochreiter_long_1997}. For network $j\in\{\mu,1,2\}$,
%
\begin{equation}
h_t^{j}=\operatorname{LN}\!\left(\operatorname{LSTM}_{j}\!\left(\operatorname{LN}(S_t);\theta_j^L\right)_{L}\right),
\label{eq:lstm_encoding}
\end{equation}
%
where the subscript $L$ selects the last sequence output. The input LayerNorm standardises each 88-dimensional observation inside the network, and a second LayerNorm is applied to the 64-dimensional LSTM output. The actor retains outputs at all sequence positions internally but acts only from the final one; each critic likewise uses only its last recurrent representation. The one-layer, 64-unit design replaces the earlier two-layer, 128-unit specification.

\subsubsection{TD3 Algorithm}
\label{subsubsec:3_td3_algorithm}

Let $\mathcal{P}$ denote the constrained long-short map and let the deterministic policy be $a_t=\mathcal{P}(\mu(S_t;\theta^\mu))$. TD3 maintains online critics $Q_1,Q_2$, an actor $\mu$, and a target copy of each network. For replay transition $(S_t,a_t,r_t,S_{t+1},d_t)$, target-policy smoothing is applied to the raw target-actor output before portfolio projection:
%
\begin{equation}
\widetilde a_{t+1}=\mathcal{P}\!\left(\mu'(S_{t+1})+\operatorname{clip}(\epsilon,-c,c)\right), \quad
\epsilon\sim\mathcal{N}(0,\sigma_{\mathrm{targ}}^2I),
\label{eq:td_target1}
\end{equation}
%  
\begin{equation}
y_t=r_t+(1-d_t)\gamma\min_{k\in\{1,2\}}Q'_k(S_{t+1},\widetilde a_{t+1}),
\label{eq:td_target2}
\end{equation}
%
with $\sigma_{\mathrm{targ}}=0.10$, $c=0.25$, and $\gamma=1$. Both critics are updated on every gradient step using
%
\begin{equation}
\mathcal{L}_Q=\mathbb{E}_{\mathcal{D}}\left[(Q_1(S_t,a_t)-y_t)^2+(Q_2(S_t,a_t)-y_t)^2\right],
\label{eq:critic_loss}
\end{equation}
%
where $\mathcal{D}$ is a uniform replay buffer with capacity 100,000 \cite{mnih_human-level_2015}.

The actor is updated once every $d=2$ critic updates using critic 1 and the differentiable action map. In addition to the deterministic policy objective \cite{silver_deterministic_2014}, the implementation uses a small allocation-entropy reward and a soft penalty on excessive leverage-gate use:
%
\begin{equation}
\mathcal{L}_{\mu}=-\mathbb{E}_{\mathcal{D}}[Q_1(S_t,\mathcal{P}(\mu(S_t)))]
-\beta_H\mathbb{E}_{\mathcal{D}}[H_t]
+\beta_g\mathbb{E}_{\mathcal{D}}[(g_t-g_0)_+^2],
\label{eq:actor_objective}
\end{equation}
%
where $\beta_H=0.005$, $\beta_g=0.25$, and $g_0=0.80$. The entropy term averages diversification within the long and active short books; it regularises the deterministic allocations and does not make TD3 a stochastic policy. After each delayed actor update, all target networks receive the Polyak update
%
\begin{equation}
\theta^{(\cdot)'}\leftarrow\tau\theta^{(\cdot)}+(1-\tau)\theta^{(\cdot)'},\qquad \tau=0.005.
\label{eq:soft_update}
\end{equation}

\subsubsection{Network Architectures}
\label{subsubsec:3_network_architectures}

The actor applies a 64-unit fully connected ReLU layer to its recurrent output and then produces $N+1$ unconstrained outputs:
%
\begin{equation}
z_t=W_2^\mu\operatorname{ReLU}(W_1^\mu h_t^\mu+b_1^\mu)+b_2^\mu,
\qquad a_t=\mathcal{P}(z_t).
\label{eq:actor_network}
\end{equation}
%
The final leverage-head weights are initialised to zero and its bias is set so that $g_t=0.10$ before learning, providing a conservative approximately unlevered starting policy. Each critic has its own LayerNorm--LSTM stack followed by one 64-unit ReLU layer:
%
\begin{equation}
Q_k(S_t,a_t)=W_{2,k}^Q\operatorname{ReLU}\!\left(W_{1,k}^Q[h_t^k;a_t]+b_{1,k}^Q\right)+b_{2,k}^Q,
\qquad k\in\{1,2\}.
\label{eq:critic_network}
\end{equation}
%
The actor and critics use separate Adam optimisers \cite{kingma_adam_2015}; gradient norms are clipped at one. Publication runs use deterministic PyTorch algorithms and full float32 arithmetic by default, with separate gradient scalers retained only for the optional mixed-precision mode.

\subsubsection{Exploration in Raw Action Space}
\label{subsubsec:3_exploration}

During the first 1,000 pre-training decisions, raw $N+1$ action vectors are sampled from a standard normal distribution to populate replay memory with broad action support. Thereafter, independent Gaussian exploration is added to the actor's raw asset-score and leverage outputs before the constrained projection:
%
\begin{equation}
a_t=\mathcal{P}\!\left(\mu(S_t)+\sigma_e\xi_t\right),
\qquad \xi_t\sim\mathcal{N}(0,I).
\label{eq:exploration_action}
\end{equation}
%
The episode-level scale is $0.15\times0.998^e$ during synthetic pre-training and $0.05\times0.995^e$ during historical fine-tuning. Adding noise before projection preserves every hard portfolio constraint. The action counter is stored in the checkpoint, so the initial random-action phase cannot restart when the pretrained agent is loaded for fine-tuning.

Algorithm \ref{alg:lstm_td3} summarises the complete LSTM-TD3 procedure, including the two-stage training strategy described next.


\subsection{Two-Stage Training Strategy}
\label{subsec:3_two_stage_training}

A central innovation of our framework is the separation of synthetic pre-training from strictly pre-holdout historical adaptation. The first stage trains TD3 on generated 24-month episodes, while the second stage starts from the resulting checkpoint and performs one controlled pass over realised historical trajectories. Both stages use the same recurrent observation, constrained action map, cost model, CRRA--CVaR reward, and one gradient update per environment step. The final 24 monthly holding periods are excluded from marginal and vine fitting, episode generation, fine-tuning, and all training diagnostics.

\subsubsection{Stage 1: Pre-Training on Synthetic Data}
\label{subsubsec:3_stage1_training}

The current configuration generates $E_{\mathrm{synth}}=1{,}000$ synthetic episodes from the training-only neural dynamic D-vine. Each episode contains 30 causal burn-in months, followed by 24 decision months. For every decision, one synthetic draw is designated as the realised return and 512 additional joint draws are retained for the CVaR calculation. The trainer verifies bundle metadata---including synthetic realised-return provenance, dynamic-vine structure, sequence length, training cutoff, and the reserved 24-period evaluation horizon---before the first update. It then uses every generated episode exactly once:

\begin{equation}
\theta_{\mathrm{pre}}=\operatorname{TD3}\!\left(\mathcal{D}_{\mathrm{synth}};E_{\mathrm{synth}}=1000,B=128,\sigma_e=0.15\times0.998^e\right).
\label{eq:pretrain_objective}
\end{equation}

Pre-training uses actor and critic learning rates $3\times10^{-5}$ and $10^{-4}$, respectively. Once training finishes, the actor is evaluated deterministically, without exploration noise, on a fixed held-in window of 100 synthetic episodes. Fine-tuning fails closed unless preregistered checks on leverage-gate use, gross-cap frequency, turnover, gate-to-gross identity, true position-limit violations, normalised allocation entropy, and effective number of positions all pass. Contact with a valid hard position cap is recorded as a warning rather than treated as a constraint violation. This gate tests the learned behaviour before any historical adaptation and does not select a policy using historical returns.

\subsubsection{Stage 2: Fine-Tuning on Real Data}
\label{subsubsec:3_stage2_training}

Fine-tuning uses realised historical monthly returns only. The configured bundle contains 61 distinct, overlapping 24-month trajectories from the training prefix; the final OOS block does not appear in any of them. To diagnose adaptation without target overlap, the last 24 months of the training prefix form one chronological validation trajectory, the 23 earlier trajectories that overlap those targets are purged, and the first 37 trajectories form the diagnostic fit. The number of passes is preregistered as one, so the validation reward is logged but cannot adaptively choose training duration.

After this diagnostic, the candidate and its replay memory are discarded. A fresh agent is loaded from the same pretrained checkpoint and refitted once on all 61 permissible historical trajectories in a deterministic seeded permutation:

\begin{equation}
\theta_{\mathrm{final}}=\operatorname{TD3}\!\left(\mathcal{D}_{\mathrm{hist}};\theta_{\mathrm{pre}},E_{\mathrm{hist}}=61,B=64,\sigma_e=0.05\times0.995^e\right).
\label{eq:finetune_objective}
\end{equation}

The actor learning rate is reduced to $2\times10^{-5}$ and the critic rate remains $10^{-4}$. This is a one-pass refit, not a 2,000-episode continuation and not an early-stopping search. The replay buffer is newly initialised for the final refit, whereas optimiser state, the action counter, and architecture metadata are restored from the pretrained checkpoint; all three target networks are then re-synchronised to their loaded online counterparts. Consequently, the random warm-up does not restart and the final policy uses all permissible historical trajectories without consulting the locked evaluation horizon.

\begin{figure}[H]
\centering
\input{figures/tikz/figure_m02_training_strategy.tex}
\caption{Implemented Two-Stage Training Strategy for Vine-Augmented LSTM-TD3}
\label{fig:two_stage_training}
\end{figure}



\subsubsection{Curriculum Learning (Optional)}
\label{subsubsec:3_curriculum_learning}

While our main implementation uses a fixed pre-training stage followed by fine-tuning, the framework can be extended with curriculum learning, where the agent is trained on progressively harder tasks \cite{bengio_curriculum_2009}. Starting with easy synthetic scenarios (low volatility, moderate correlations) and progressing to harder scenarios (high volatility, extreme correlations) can improve sample efficiency and final policy performance. The training data at each curriculum stage is:

\begin{equation}
\mathcal{D}_k = \begin{cases}
\mathcal{D}_{\text{synth}}^{(\text{easy})}, & k = 1 \\
\mathcal{D}_{\text{synth}}^{(\text{medium})}, & k = 2 \\
\mathcal{D}_{\text{synth}}^{(\text{hard})}, & k = 3 \\
\mathcal{D}_{\text{real}}, & k = 4
\end{cases}
\label{eq:curriculum}
\end{equation}

This approach gradually introduces complexity, allowing the agent to build robust policies before facing the full challenge of real market data. While not employed in our main experiments, this extension demonstrates the flexibility of the framework.

Algorithm \ref{alg:lstm_td3} presents the complete training procedure, including both pre-training and fine-tuning stages.

\begin{algorithm}[H]
\caption{LSTM-TD3 Training Procedure}
\label{alg:lstm_td3}
\begin{algorithmic}[1]
\STATE \textbf{Input:} Synthetic episodes $\mathcal{D}_{\text{synth}}$, historical training-prefix episodes $\mathcal{D}_{\text{hist}}$, hyperparameters
\STATE \textbf{Output:} Trained policy $\mu(s|\theta^\mu)$
\STATE
\STATE Validate episode metadata, causal burn-in histories, and locked 24-period OOS exclusion
\STATE Initialize recurrent actor $\mu$, twin recurrent critics $Q_1,Q_2$, and their target copies
\STATE Initialize replay buffer $\mathcal{B}$ with capacity $N_B$
\STATE
\STATE \textbf{Stage 1: Pre-training on Synthetic Data}
\FOR{each of 1,000 synthetic episodes}
    \STATE Reset from its 30-month causal burn-in sequence $S_0$
    \FOR{$t=1$ to $24$}
        \STATE Select raw random action during the first 1,000 decisions; otherwise use $\mu(S_t)+\sigma_e\xi_t$
        \STATE Project raw outputs to $a_t=\mathcal{P}(z_t)$; execute and observe $(r_t,S_{t+1},d_t)$
        \STATE Store $(S_t,a_t,r_t,S_{t+1},d_t)$ in $\mathcal{B}$ and sample a uniform mini-batch when $|\mathcal{B}|\geq128$
        \STATE Form smoothed target actions and $y_i=r_i+(1-d_i)\min(Q'_1,Q'_2)$
        \STATE Update both critics; clip each gradient norm at 1
        \IF{critic update count is divisible by 2}
            \STATE Update the actor with $-Q_1$, entropy regularisation, and the leverage penalty
            \STATE Soft-update the actor and both critic targets with $\tau=0.005$
        \ENDIF
    \ENDFOR
\ENDFOR
\STATE Save the pretrained checkpoint and evaluate the deterministic actor on 100 held-in synthetic episodes
\STATE Abort unless every preregistered behavioural gate passes
\STATE
\STATE \textbf{Stage 2: Fine-tuning on Real Data}
\STATE Load a fresh copy of the pretrained agent; fit once on 37 target-disjoint trajectories and log one purged validation trajectory
\STATE Discard the diagnostic candidate and replay buffer
\STATE Reload the pretrained agent and train once on all 61 historical trajectories in a seeded permutation, using $B=64$
\STATE Save the full checkpoint, diagnostics, episode schedule, and validation log
\RETURN $\mu(s|\theta^\mu)$
\end{algorithmic}
\end{algorithm}











\subsection{Benchmark Models}
\label{subsec:3_benchmark_models}

The locked main evaluation contains six causal, non-RL benchmarks: equal weight, constrained shrinkage mean--variance, DCC--GARCH, static vine, rolling vine, and dynamic NN-vine optimisation. Every strategy receives the same dates, asset order, realised returns, transaction and financing costs, and long--short mandate. A benchmark may use only observations dated no later than its decision date; solver failure is recorded and never replaced by a silent fallback.

\subsubsection{Classical Models}
\label{subsubsec:3_classical_models}

Equal weight provides a transparent passive reference. The optimised classical benchmarks are shrinkage mean--variance, which regularises the covariance matrix before solving the Markowitz objective \cite{markowitz_portfolio_1952}, and the Dynamic Conditional Correlation model of Engle \cite{engle_dynamic_2002}. At decision date $t$, the generic problem is

\begin{equation}
w^* = \arg\min_{w \in \mathcal{A}} \frac{1}{2} w^\top \Sigma w - \lambda w^\top \mu,
\label{eq:markowitz_constrained}
\end{equation}

where $\Sigma_t$ is the causal covariance estimate, $\mu_t$ the causal mean estimate, and the feasible set $\mathcal A$ is the same net-one, gross-1.5, $[-0.20,0.60]$ set as the RL policy. Shrinkage stabilises $\Sigma_t$ in the short sample. DCC--GARCH instead combines univariate conditional volatilities with a time-varying correlation matrix and solves

\begin{equation}
w_t^* = \arg\min_{w \in \mathcal{A}} \frac{1}{2} w^\top \Sigma_t w - \lambda w^\top \hat{\mu}_t.
\end{equation}

This captures time-varying linear dependence but, like all elliptical models, fails to accommodate non-linear or asymmetric tail dependencies.

\subsubsection{Vine-Based Models}
\label{subsubsec:3_vine_based_models}

The three vine benchmarks isolate the value of nonlinear dependence from the value of learned sequential control. The \textit{Static D-Vine} estimates one training-prefix vine and reuses it at every decision. The \textit{Rolling D-Vine} refits on the available lookback at each date. The \textit{Dynamic NN-D-Vine} uses the neural parameterisation of Section \ref{subsubsec:3_dynamic_vine_with_NN_param} without RL. Each model generates 512 causal return scenarios and chooses a portfolio by constrained one-period CRRA expected utility:

\begin{equation}
w_t^*=\arg\max_{w\in\mathcal A}\frac{1}{512}\sum_{m=1}^{512}
U_{\gamma_U}\!\left(g_{p,t}^{(m)}(w)m_t(w)\right),
\label{eq:myopic_eu}
\end{equation}

where the same trading and financing costs are included in the scenario growth factors. The dynamic NN-vine optimiser is the most direct non-RL comparator: it shares the proposed dependence engine and scenario distribution but replaces the recurrent policy with a transparent myopic optimiser.

\subsubsection{RL Baseline}
\label{subsubsec:3_rl_baseline}

The frozen coarse \textit{no-vine LSTM-TD3} ablation retains the recurrent architecture, input dimension, long--short action map, CRRA--CVaR reward, two-stage training protocol, costs, and matched training logic. It sets both the 63-dimensional vine vector and the policy-visible scenario-CVaR component to zero while leaving their input positions present. Because that intervention removes two channels simultaneously, it is an informative systems ablation but not a clean estimate of the direct vine-state effect.

\subsubsection{Matched-Seed Causal Mechanism Analysis}
\label{subsubsec:3_causal_analysis}

To identify mechanisms more precisely, a prospective post-holdout contract defines 13 experiments on the same architecture and data pipeline. The reference is full TD3 with vine state, scenario-CVaR observation, and CVaR reward. Eight primary interventions separately remove direct vine features, the scenario-CVaR observation, both policy-visible dependence channels, the CVaR reward term, synthetic pre-training, recurrence, and historical fine-tuning; the eighth replaces neural-vine pre-training with a moving-block-bootstrap bundle. Four exploratory algorithm controls replace TD3 with DDPG, SAC, PPO, or A2C while holding the observation, reward, cost, constraint, and environment-interaction contracts fixed.

Every experiment uses the matched seeds 20261001--20261010. Seeds measure optimisation variability only; they are not treated as independent market histories. For each variant, target weights are averaged across seeds and the ensemble is rescored once with drifted pre-trade holdings. The primary outcome is the paired monthly CRRA-utility difference relative to the reference, supplemented by the annualised certainty-equivalent effect. Inference uses 9,999 circular moving-block bootstrap draws of length three. One-sided p-values for the eight component contrasts are Holm-adjusted as one family; the four algorithm comparisons form a separate two-sided exploratory family. Because the market holdout had already been evaluated when this protocol was written, all causal results are labelled post-holdout explanatory and cannot restore confirmatory status.




\subsection{Implementation Details}
\label{subsec:3_implementation}

This section describes the software and hardware environment, the hyperparameter configuration, and the key implementation choices necessary for reproducibility of the empirical results.

\subsubsection{Software and Hardware}
\label{subsubsec:3_software_hardware}

Experiments are executed on a Linux HPC cluster with AMD EPYC CPUs and NVIDIA A100-SXM4-80GB GPUs under Conda/Miniforge environments. Data handling, vine estimation, scenario generation, and the portfolio environment are implemented in \textsf{R}; neural policies are implemented in \textsf{Python}/PyTorch. The exact executable paths, package inventories, CUDA runtime, source hashes, and data hashes are stored with each frozen release rather than inferred from the interactive shell. Training uses the GPU PyTorch environment, whereas evaluation launches policy inference in a separate CPU PyTorch process. This process isolation prevents the R Lantern and Python libtorch libraries from loading incompatible symbols into one address space. The implementation uses the following resource controls:
\begin{itemize}
    \item \textbf{Parallel simulation}: The \texttt{rvinecopulib} simulator requests the configured physical-core budget for synthetic episodes and scenarios, capped at the cores detected by the worker.
    \item \textbf{Pre-computed episodes}: The current configuration stores 1,000 synthetic 24-step episodes and 61 realised historical 24-step trajectories in portable \texttt{qs}/R bundles. Each episode also carries a 30-step causal burn-in, dynamic vine states, and the scenario matrices required for CVaR, avoiding vine re-estimation inside the PyTorch update loop.
    \item \textbf{GPU execution}: One training process uses one CUDA device. Independent seed replicates are sharded across the available two or four GPUs; a single agent is not data-parallel. Deterministic PyTorch algorithms and float32 arithmetic are the publication defaults, with automatic mixed precision disabled.
    \item \textbf{Update frequency}: The trainer performs one twin-critic update per environment transition and one actor/target update per two critic updates. The former three-updates-per-step schedule is no longer used.
\end{itemize}
The complete software stack, configuration, source hashes, data hashes, checkpoint schema, and diagnostic files are captured in the run artifacts. Training and evaluation runtimes are reported from those manifests; timings from the superseded DDPG schedule are not used.


\begin{table}[htbp]
\centering
\caption{Software and Hardware Specifications}
\label{tab:hardware}
\begin{tabular}{ll}
\toprule
\textbf{Component} & \textbf{Specification} \\
\midrule
CPU & AMD EPYC; worker allocation recorded per run (up to 128 cores) \\
GPU & NVIDIA A100-SXM4-80GB; one GPU per training worker \\
RAM & HPC allocation recorded in the run manifest \\
Operating System & Linux \\
Environment Manager & Conda (Miniforge3) \\
R Stack & \texttt{rvinecopulib}, \texttt{rugarch}, \texttt{torch}/Lantern, \texttt{xts}, \texttt{qs} \\
Training Python & Frozen GPU environment; PyTorch 2.10.0+cu128 in the causal sweep \\
Evaluation Python & Isolated CPU environment; PyTorch 2.7.1+cpu \\
Inter-process Interface & Immutable file IPC for policy observations and actions \\
\bottomrule
\end{tabular}
\end{table}

\subsubsection{Hyperparameter Configuration}
\label{subsubsec:3_hyperparameter_config}

Table \ref{tab:hyperparameters} reports the values frozen in the current master configuration. The recurrent lookback is 30 monthly observations, not 30 trading days. Because dense utility increments telescope to terminal-wealth utility, the TD discount is fixed at one. The environment uses CRRA coefficient 2, CVaR penalty 0.05, proportional turnover cost 0.001, and 512 scenarios for the 95\% CVaR estimate. Hyperparameters are not selected on the final 24-period OOS path.

\begin{table}[H]
\centering
\caption{Current LSTM-TD3 and Training Hyperparameters}
\label{tab:hyperparameters}
\resizebox{\textwidth}{!}{%
\begin{tabular}{l c c l c c}
\toprule
\textbf{Parameter} & \textbf{Symbol} & \textbf{Value} & \textbf{Parameter} & \textbf{Symbol} & \textbf{Value} \\
\midrule
LSTM sequence length & $L$ & 30 months & Batch size (pre/fine) & $B$ & 128 / 64 \\
LSTM hidden units & — & 64 & Replay buffer capacity & $N_B$ & 100,000 \\
LSTM layers & — & 1 & Discount factor & $\gamma$ & 1.0 \\
Actor learning rate (pre/fine) & $\alpha_\mu$ & $3\times10^{-5}$ / $2\times10^{-5}$ & Target update rate & $\tau$ & 0.005 \\
Critic learning rate (pre/fine) & $\alpha_Q$ & $10^{-4}$ / $10^{-4}$ & Policy delay & $d$ & 2 \\
Target noise / clip & $\sigma_{\rm targ},c$ & 0.10 / 0.25 & Episodes (pre/fine) & — & 1,000 / 61 \\
Exploration scale (pre/fine) & $\sigma_e$ & 0.15 / 0.05 & Exploration decay & — & 0.998 / 0.995 \\
CVaR penalty / scenarios & $\lambda,M$ & 0.05 / 512 & Transaction cost & $\kappa$ & 0.001 \\
Utility risk aversion & $\gamma_U$ & 2 & Gross / net exposure & — & 1.5 / 1.0 \\
Maximum long / short & — & 0.60 / 0.20 & Entropy coefficient & $\beta_H$ & 0.005 \\
Leverage target / penalty & $g_0,\beta_g$ & 0.80 / 0.25 & Gradient clip norm & — & 1.0 \\
\bottomrule
\end{tabular}%
}
\end{table}

Exploration uses independent Gaussian noise in the actor's raw output space, following the stage-specific scales and decay rates in Table \ref{tab:hyperparameters}. Target-policy smoothing is a separate TD3 mechanism and uses its own fixed standard deviation and clipping threshold.

\subsubsection{Training Procedure}
\label{subsubsec:3_training_procedure}

The training procedure follows Section \ref{subsec:3_two_stage_training}: 1,000 synthetic episodes are used once for pre-training, a deterministic 100-episode policy check must pass, and one historical pass is then performed. The diagnostic fine-tuning fit and the final all-history refit each begin from the same pretrained checkpoint and use separate replay buffers. The historical validation path is diagnostic because the configuration fixes the pass count at one; there is no 100-epoch early-stopping rule.

At each environment step, the trainer appends one transition and, once the replay buffer contains a full batch, samples uniformly. Batch size is 128 during pre-training and 64 during fine-tuning. Adam updates both critics on every step and the actor on alternating updates; all recurrent and fully connected gradients are clipped at norm one. Non-finite actions, rewards, actor losses, or critic losses terminate the run rather than being silently skipped.

Monitoring is artifact-based rather than a subjective reward-plateau test. Episode CSV files record reward, terminal wealth, turnover, CVaR, gross exposure, short notional, leverage-gate behaviour, position-cap contact, and allocation entropy. Update CSV files record both Q estimates, their target and disagreement, losses, gradient norms, actor regularisers, and precision scales. Checkpoint metadata locks the TD3 architecture, recurrent dimensions, action schema, mandate, costs, utility mode, vine-observation mode, and schema version before evaluation can load the policy.

\subsubsection{Reproducibility and Code Availability}
\label{subsubsec:3_reproducibility}

The master configuration uses seed 20260741 for the first full-policy run and 20260741--20260760 for the frozen 20-policy ensemble. Seeds are applied to R, NumPy, Python, and PyTorch. Deterministic algorithms, code and data hashes, the exact configuration, checkpoint architecture metadata, and a source snapshot are written for every run. Seed replications quantify optimisation instability; paired time-series inference is performed over realised holding periods and never treats seeds as additional market observations. The locked main batch, frozen releases, operational-retry logs, and post-holdout evidence classes are retained so that unsuccessful execution attempts and later explanatory analyses remain distinguishable from the original evaluation.

\section{Data and Empirical Design}
\label{sec:data}

\subsection{Asset Universe and Return Construction}
\label{subsec:4_asset_selection}

The empirical panel contains seven adjusted market-level series spanning US equities, Chinese equities, and gold. The fixed column order is S\&P 500, NASDAQ, Dow Jones, SSE 50, Dividend Index, ChiNext, and Gold. The universe was selected before the locked evaluation and is held constant in model fitting, checkpoint metadata, benchmark generation, and scoring. Table \ref{tab:assets} summarises the economic role of each series.

\begin{table}[H]
\centering
\caption{Seven-Asset Portfolio Universe}
\label{tab:assets}
\begin{tabular}{l l l}
\toprule
\textbf{Code} & \textbf{Market exposure} & \textbf{Portfolio role} \\
\midrule
SP500 & Large-cap US equities & Broad US equity factor \\
NASDAQ & Technology-heavy US equities & Growth and technology factor \\
DOW & US blue-chip equities & Mature large-cap factor \\
SSE50 & Large-cap mainland Chinese equities & China large-cap factor \\
DIVIDEND & Chinese dividend equities & China value/income factor \\
CHINEXT & Chinese growth equities & China growth factor \\
GOLD & Gold adjusted-level series & Defensive real-asset factor \\
\bottomrule
\end{tabular}
\end{table}

The immutable input contains 2,946 synchronized adjusted-level observations from 16 December 2013 to 6 July 2026. Dates are required to be valid, unique, and strictly increasing; all seven levels must be finite and strictly positive. Daily log returns are constructed as

\begin{equation}
r_{i,t}=\log P_{i,t}-\log P_{i,t-1},
\label{eq:data_log_returns}
\end{equation}

which yields 2,945 complete common-date return rows. The loader neither reorders dates nor interpolates, winsorises, forward-fills, or silently drops an asset. Monthly portfolio returns are compounded from the daily log returns over $(d_t,d_{t+1}]$, where $d_t$ is the final common trading date available for the decision month. This common calendar is essential because US and Chinese market holidays differ.

\subsection{Descriptive Statistics}
\label{subsec:4_descriptive_statistics}

Table \ref{tab:descriptive} reports full-panel daily log-return summaries. These figures describe the realised sample; they are not inputs estimated with access to future data. Marginal and copula parameters used by the agent are fitted only on the training prefix described below.

\begin{table}[H]
\centering
\caption{Daily Log-Return Descriptive Statistics, December 2013--July 2026}
\label{tab:descriptive}
\resizebox{\textwidth}{!}{%
\begin{tabular}{l r r r r r r}
\toprule
\textbf{Asset} & \textbf{Mean (\%)} & \textbf{SD (\%)} & \textbf{Skewness} & \textbf{Excess kurtosis} & \textbf{Minimum (\%)} & \textbf{Maximum (\%)} \\
\midrule
S\&P 500 & 0.0489 & 1.1112 & -0.631 & 16.145 & -12.765 & 9.089 \\
NASDAQ & 0.0635 & 1.3433 & -0.406 & 8.848 & -13.141 & 11.478 \\
Dow Jones & 0.0410 & 1.0892 & -0.819 & 22.081 & -13.842 & 10.764 \\
SSE 50 & 0.0218 & 1.3923 & -0.217 & 8.806 & -10.517 & 11.150 \\
Dividend Index & 0.0187 & 1.3342 & -0.944 & 13.118 & -10.540 & 9.538 \\
ChiNext & 0.0391 & 2.1472 & -0.062 & 8.901 & -17.770 & 18.232 \\
Gold & 0.0448 & 0.9826 & -0.774 & 14.151 & -10.533 & 7.212 \\
\bottomrule
\end{tabular}%
}
\caption*{\footnotesize Notes: Statistics use all 2,945 synchronized daily log returns and are descriptive. Excess kurtosis is relative to the Gaussian value of zero.}
\end{table}

All seven series are heavy-tailed and six are materially negatively skewed. ChiNext has the greatest daily volatility and the widest observed range; Gold has the lowest volatility. These properties motivate asset-specific conditional volatility filters and a tail-dependent copula rather than a homoskedastic multivariate Gaussian model.

\subsection{Cross-Asset Dependence}
\label{subsec:4_correlation}

Table \ref{tab:correlation} gives the full-panel Pearson correlation matrix. Correlations within the three US indices are high, especially S\&P 500--Dow Jones (0.948) and S\&P 500--NASDAQ (0.945). Chinese equity dependence is also substantial but lower, with SSE 50--Dividend at 0.757 and SSE 50--ChiNext at 0.566. Cross-region equity correlations lie near 0.10--0.16, while Gold is close to uncorrelated with every equity series. These unconditional values conceal conditional and tail dependence, which the dynamic D-vine is designed to represent.

\begin{table}[H]
\centering
\caption{Unconditional Pearson Correlation Matrix}
\label{tab:correlation}
\begin{tabular}{l c c c c c c c}
\toprule
 & \textbf{SP500} & \textbf{NASDAQ} & \textbf{DOW} & \textbf{SSE50} & \textbf{DIV} & \textbf{CNXT} & \textbf{GOLD} \\
\midrule
\textbf{SP500} & 1.000 & 0.945 & 0.948 & 0.155 & 0.109 & 0.142 & 0.016 \\
\textbf{NASDAQ} & 0.945 & 1.000 & 0.826 & 0.156 & 0.097 & 0.149 & 0.012 \\
\textbf{DOW} & 0.948 & 0.826 & 1.000 & 0.146 & 0.110 & 0.122 & 0.019 \\
\textbf{SSE50} & 0.155 & 0.156 & 0.146 & 1.000 & 0.757 & 0.566 & 0.052 \\
\textbf{DIV} & 0.109 & 0.097 & 0.110 & 0.757 & 1.000 & 0.499 & 0.046 \\
\textbf{CNXT} & 0.142 & 0.149 & 0.122 & 0.566 & 0.499 & 1.000 & 0.038 \\
\textbf{GOLD} & 0.016 & 0.012 & 0.019 & 0.052 & 0.046 & 0.038 & 1.000 \\
\bottomrule
\end{tabular}
\caption*{\footnotesize Notes: Pearson correlations of synchronized daily log returns over the full descriptive sample.}
\end{table}

\subsection{Marginal and Synthetic-Data Diagnostics}
\label{subsec:4_stat_autocorrelation_tests}

The original draft relied on generic full-sample ADF and Ljung--Box tables. The implemented protocol instead gates the objects that enter the copula. For every asset, candidate marginal models are estimated on the training prefix and the standardised residuals and squared residuals are tested at lags 5, 10, and 20. The model proceeds only if the selected marginal has zero-mean, approximately unit-variance residuals and passes the preregistered 1\% residual-autocorrelation gates. This is stricter and more directly relevant than using a unit-root test on returns as evidence for a particular GARCH form.

Synthetic validation is also multi-dimensional. Marginal summaries compare mean, dispersion, quantiles, skewness, kurtosis, VaR, and CVaR. Cross-sectional diagnostics compare all 21 Pearson correlations and lower-tail co-exceedance counts with finite-sample compatibility intervals. Temporal diagnostics compare lagged returns and squared returns. In the frozen generator run, marginal fidelity passed for all seven assets, the strict correlation rule passed for 20 of 21 pairs, and lower-tail co-exceedance was statistically compatible for all 21 pairs. The single strict-correlation miss is retained and reported rather than tuned away; passing diagnostics supports fidelity to the measured training sample, not proof that synthetic paths are indistinguishable from future markets.

\subsection{Data Partitioning and Locked Evaluation}
\label{subsec:4_data_partitioning}

Monthly periods are built first and only then split. The final 24 decisions, from 31 July 2024 through 30 June 2026 with holding ends through 6 July 2026, form the locked evaluation block. All marginal, vine, synthetic-data, RL, and benchmark-selection operations are restricted to information dated no later than the corresponding training or decision cutoff. The training-prefix holding periods end no later than 31 July 2024, so the first evaluation return is never present in training.

Synthetic pre-training uses 1,000 newly simulated 24-step episodes and no realised historical outcomes. Historical fine-tuning uses all 61 permissible overlapping 24-month trajectories from the prefix. A diagnostic fit uses 37 early trajectories and a chronologically later 24-month validation path after purging the 23 trajectories whose targets overlap validation; because the number of passes is fixed at one, validation is logged but does not select duration. The final agent is reloaded from the same pretrained checkpoint and fitted once on all 61 trajectories.

The locked calendar contains 24 planned holding periods. The preregistered primary scope requires at least 15 joint trading observations and excludes two rows: 30 January--27 February 2026 has 14 common observations, and 30 June--6 July 2026 has three. Primary economic and inferential results therefore use 22 complete periods. The 24-row calculation is reported only as shortened-period sensitivity evidence and does not replace the declared primary sample.

\subsection{Remarks on Data Quality and Limitations}
\label{subsec:4_remarks}

Several limitations remain. First, the panel contains broad adjusted-level series rather than directly tradable total-return ETFs, so fees, tracking error, dividends, tax treatment, and currency conversion are not fully represented. Second, synchronized common dates avoid interpolation but can leave relatively short cross-market holding periods; the 15-observation completeness rule is conservative and should be frozen separately for future panels. Third, the sample contains only one realised market path and 22 complete primary test periods. Monthly 5\% CVaR is therefore based on one observed tail month and is descriptive, not a reliable tail-risk superiority test. Finally, the present seven-asset universe is modest. External validity requires a separately frozen investable multi-asset panel and non-overlapping walk-forward windows, with source licences, total-return definitions, and immutable data hashes included in the replication package.

\section{Results and Discussion}
\label{sec:results}

This section reports the frozen main evaluation before turning to post-holdout explanatory analyses. The distinction is important. The main batch was specified before successful scoring of the locked path and contains 20 full policies and six benchmarks. The no-vine, pretrained-only, ensemble-mechanism, and causal analyses were conducted after that path had been accessed and therefore explain mechanisms but cannot provide a second independent confirmation.

\subsection{Evaluation Contract and Performance Measures}
\label{subsec:4_evaluation_contract}

Every strategy is evaluated on the same realised seven-asset return matrix, calendar, asset order, initial wealth of 100,000, and cost convention. The RL deployment is the preregistered arithmetic mean of the 20 seed-specific target weights; returns are never averaged across agents. The common scorer first verifies the net, gross, long, and short constraints and then recomputes portfolio growth and costs from the ensemble weights.

For $n$ complete holding periods spanning $Y$ elapsed years, total return and compound annual growth are

\begin{equation}
R_{\mathrm{tot}}=\frac{V_n}{V_0}-1,
\qquad
R_{\mathrm{CAGR}}=\left(\frac{V_n}{V_0}\right)^{1/Y}-1.
\label{eq:annualized_return}
\end{equation}

Volatility is annualised from monthly net returns with factor $\sqrt{12}$. Sharpe and Sortino ratios use a zero risk-free rate; maximum drawdown is computed from the cumulative wealth path. Omega is the ratio of gains above zero to losses below zero. The annualised CRRA certainty-equivalent return, the primary economic measure, is obtained from mean monthly utility $\bar U$ as

\begin{equation}
R_{\mathrm{CE}}=\left[1+(1-\gamma_U)\bar U\right]^{12/(1-\gamma_U)}-1,
\qquad \gamma_U=2.
\label{eq:realized_utility}
\end{equation}

All annualisation uses the recorded holding calendar rather than the former daily factor of 252. With only 22 complete months, realised 5\% VaR and CVaR depend on one tail event and are treated as descriptive diagnostics.

\subsection{Locked Out-of-Sample Economic Performance}
\label{subsec:4_outofsample}

Table \ref{tab:outofsample} reports the primary 22-complete-period results. The NN-vine LSTM-TD3 ensemble finishes with 159,668.87, a total return of 59.67\%. It has the lowest observed annual volatility (13.19\%) and the highest observed Sharpe ratio (2.018). It does not, however, maximise wealth or the primary utility measure. The static vine has the highest total return (64.50\%) and certainty-equivalent return (29.55\%); the TD3 ensemble ranks third on both measures, behind the static and dynamic NN-vine optimisers.

\begin{table}[H]
\centering
\caption{Frozen Out-of-Sample Performance: 22 Complete Holding Periods}
\label{tab:outofsample}
\resizebox{\textwidth}{!}{%
\begin{tabular}{l r r r r r r r r}
\toprule
\textbf{Strategy} & \textbf{Total} & \textbf{CAGR} & \textbf{Vol.} & \textbf{Sharpe} & \textbf{Sortino} & \textbf{Max DD} & \textbf{Omega} & \textbf{CRRA CE} \\
\midrule
Equal weight & 49.16\% & 24.37\% & 13.42\% & 1.702 & \textbf{5.029} & \textbf{5.22\%} & 4.141 & 23.38\% \\
Shrinkage mean--variance & 48.69\% & 24.16\% & 14.17\% & 1.608 & 2.952 & 7.83\% & 3.363 & 22.98\% \\
DCC--GARCH & 58.79\% & 28.69\% & 13.55\% & 1.944 & 4.632 & 5.79\% & \textbf{4.965} & 27.62\% \\
Static vine & \textbf{64.50\%} & \textbf{31.19\%} & 16.88\% & 1.704 & 4.478 & 6.51\% & 3.756 & \textbf{29.55\%} \\
Rolling vine & 40.81\% & 20.52\% & 18.94\% & 1.082 & 2.033 & 9.90\% & 2.359 & 18.53\% \\
Dynamic NN-vine, no RL & 61.44\% & 29.86\% & 16.45\% & 1.682 & 4.300 & 6.78\% & 3.760 & 28.30\% \\
\textbf{NN-vine LSTM-TD3 ensemble} & 59.67\% & 29.08\% & \textbf{13.19\%} & \textbf{2.018} & 4.642 & 6.64\% & 4.880 & 28.05\% \\
\bottomrule
\end{tabular}%
}
\caption*{\footnotesize Notes: Volatility, CAGR, and CRRA certainty-equivalent return are annualised. Bold denotes the best observed value, with lower values preferred for volatility and drawdown. These rankings are descriptive and do not imply statistical significance.}
\end{table}

Relative to equal weight, the ensemble creates 10,504.49 more terminal wealth, or 7.04\% of equal-weight terminal wealth. Its terminal-wealth differences are +10,981.78 versus shrinkage mean--variance, +876.88 versus DCC--GARCH, -4,832.60 versus the static vine, +18,857.95 versus the rolling vine, and -1,773.23 versus the dynamic NN-vine without RL. The empirical result is therefore economically competitive but not dominant.

The two incomplete periods are not silently discarded after observing performance: exclusion follows the preregistered minimum of 15 common trading observations. A 24-row shortened-period sensitivity makes the ensemble rank first, but this reversal is driven materially by the final three-day row, in which both non-RL vine optimisers lose approximately 3\% and the ensemble is nearly flat. The 24-row result is accordingly a sensitivity check, not the primary estimate.

\subsection{Paired Statistical Inference}
\label{subsec:4_inference}

Inference is based on paired monthly CRRA-utility differences, thereby preserving the fact that all strategies experience the same market periods. A circular moving-block bootstrap accounts for short-range temporal dependence. The primary alternative is that the ensemble has higher utility than equal weight; secondary comparisons use the same construction, and Holm adjustment controls multiplicity across the six benchmark contrasts.

\begin{table}[H]
\centering
\caption{Paired CRRA-Utility Effects of the TD3 Ensemble}
\label{tab:paired_inference}
\resizebox{\textwidth}{!}{%
\begin{tabular}{l r c r r}
\toprule
\textbf{Benchmark} & \textbf{Mean effect} & \textbf{95\% block-bootstrap interval} & \textbf{One-sided $p$} & \textbf{Holm $p$} \\
\midrule
Equal weight & 0.003034 & $[-0.005692,\ 0.010116]$ & 0.2347 & 0.9388 \\
Shrinkage mean--variance & 0.003303 & $[-0.001176,\ 0.008345]$ & 0.0928 & 0.5568 \\
DCC--GARCH & 0.000273 & $[-0.004230,\ 0.005428]$ & 0.4406 & 1.0000 \\
Static vine & -0.000951 & $[-0.009173,\ 0.008199]$ & 0.5639 & 1.0000 \\
Rolling vine & 0.006323 & $[-0.006399,\ 0.018376]$ & 0.1645 & 0.8225 \\
Dynamic NN-vine, no RL & -0.000158 & $[-0.009793,\ 0.010222]$ & 0.5024 & 1.0000 \\
\bottomrule
\end{tabular}
}
\caption*{\footnotesize Notes: Effects are ensemble minus benchmark mean monthly CRRA utility. All intervals include zero. Seeds are not used as market replications.}
\end{table}

The primary comparison fails to reject: the positive mean effect versus equal weight has one-sided $p=0.2347$ and Holm-adjusted $p=0.9388$. Every interval in Table \ref{tab:paired_inference} contains zero. White's reality check selects the static vine as the best observed candidate against equal weight, but its $p$-value is 0.5142. No strategy therefore has a statistically established advantage after accounting for the short path and strategy search.

The ensemble's realised 5\% VaR and CVaR losses are 5.29\% and 6.64\%, respectively. Since empirical CVaR is based on a single month, it cannot substantiate a tail-risk-dominance claim. The strongest defensible conclusion is that the ensemble occupies an attractive observed risk--return point; the data are insufficient to establish superiority.

\subsection{Seed Robustness and the Ensemble Mechanism}
\label{subsec:4_seed_robustness}

Training randomness materially affects the learned policy. Table \ref{tab:seed_robustness} compares the worst, median, and best individual seed with the mean-weight ensemble. The ensemble has a higher Sharpe ratio and lower turnover than the median seed, but it is not representative of the typical learned long--short book.

\begin{table}[H]
\centering
\caption{Out-of-Sample Dispersion Across 20 Training Seeds}
\label{tab:seed_robustness}
\begin{tabular}{l r r r r}
\toprule
\textbf{Metric} & \textbf{Worst seed} & \textbf{Median seed} & \textbf{Best seed} & \textbf{Ensemble} \\
\midrule
CAGR & 15.47\% & 25.95\% & 45.46\% & 29.08\% \\
Sharpe ratio & 0.983 & 1.628 & 2.237 & 2.018 \\
Maximum drawdown & 12.15\% & 8.10\% & 5.45\% & 6.64\% \\
Terminal wealth & 130,169.62 & 152,677.09 & 198,765.78 & 159,668.87 \\
Monthly turnover & 1.001 & 0.539 & 0.338 & 0.317 \\
\bottomrule
\end{tabular}
\end{table}

Only 11 of 20 seeds beat equal weight on CAGR and 9 beat it on Sharpe. Eight beat the static vine on CAGR, while nine beat either the static or dynamic NN-vine on Sharpe. No individual seed beats equal weight on maximum drawdown or empirical CVaR. These fractions reject a seed-robust superiority narrative.

Cross-seed averaging changes the economic exposure. The mean individual policy has gross exposure 1.4023, short notional 20.12\%, monthly turnover 0.5747, transaction costs of 5.747 basis points per month, and financing costs of 5.029 basis points. The ensemble has corresponding values 1.0221, 1.11\%, 0.3170, 3.170 basis points, and 0.277 basis points. Approximately 94.5\% of incremental gross and short exposure cancels because seeds take opposing positions. The high ensemble Sharpe is thus partly an ensemble diversification, de-leveraging, and cost-compression effect, not simply evidence that a typical agent learned a stable leveraged policy.

The ensemble's average allocations are approximately 33.1\% Gold, 26.9\% NASDAQ, 14.1\% Dow Jones, 13.1\% S\&P 500, 7.8\% ChiNext, 2.9\% Dividend, and 2.3\% SSE 50. This concentration benefited from the realised Gold and US growth performance in the locked path and should not be interpreted as evidence of crisis robustness.

\subsection{Implementation, Constraints, and Numerical Integrity}
\label{subsec:4_implementation}

All 648 strategy-period rows satisfy the declared constraints at tolerance $10^{-6}$. The maximum net-exposure error is $9.92\times10^{-9}$, the largest gross exposure is 1.5000000098, the largest long weight is 0.60, and the smallest short weight is -0.20. The TD3 ensemble's mean and maximum gross exposure are 1.022 and 1.121. Its implementation drag lowers total return by 1.215 percentage points relative to its pre-cost 60.884\% result.

The successful frozen main scorer used target-to-target turnover, because it preceded the later drift-aware correction described in Section \ref{subsubsec:3_reward_function}. A deterministic post-hoc reconciliation using drifted pre-trade holdings raises ensemble turnover from 0.3170 to 0.3303 and lowers total return by only 0.047 percentage points; rankings and substantive conclusions are unchanged. The frozen main table is not overwritten. All causal and future protocols use the corrected drift-aware convention with actual-calendar-day financing.

There is no silent benchmark fallback and all generated weights are finite and feasible. Nevertheless, 12 of 120 optimised benchmark decisions---four each for the static, rolling, and dynamic NN-vine strategies---terminate with NLopt code 5 at the frozen 2,000-evaluation budget. This code means maximum evaluations were reached, not that an optimality tolerance was met. The portfolios are feasible, but numerical optimality is not proven. The original result is retained; subsequent protocols fail closed on code 5 and require preregistered multi-start convergence diagnostics before a new holdout is accessed.

The locked batch itself required 319.7 seconds for six benchmark paths, approximately 11 seconds for each isolated RL-policy replay, 8.9 seconds for common scoring, and less than one second for realised-panel construction. Policy inference was isolated in a CPU PyTorch subprocess to avoid shared-library collisions with R Lantern. Failed operational attempts and the successful release are preserved with immutable logs; overlapping weights and realised panels are byte-identical, but the paper does not describe the successful run as first-ever untouched access.

\subsection{Post-Holdout Explanatory Ablations}
\label{subsec:4_ablation}

Table \ref{tab:ablation} reports two frozen same-sample ablations. The no-vine ensemble jointly masks the direct vine state and policy-visible scenario-CVaR channel. The pretrained-only ensemble omits historical fine-tuning. Because both were evaluated after the main holdout had been consumed, no confirmatory test is performed and the results are descriptive mechanism evidence.

\begin{table}[H]
\centering
\caption{Frozen Post-Holdout Explanatory Ablations: 22 Complete Periods}
\label{tab:ablation}
\resizebox{\textwidth}{!}{%
\begin{tabular}{l r r r r r r r r}
\toprule
\textbf{Variant} & \textbf{Total} & \textbf{CAGR} & \textbf{Vol.} & \textbf{Sharpe} & \textbf{Sortino} & \textbf{Max DD} & \textbf{CRRA CE} & \textbf{Turnover} \\
\midrule
Full NN-vine TD3 ensemble & 59.67\% & 29.08\% & 13.19\% & 2.018 & 4.642 & 6.64\% & 28.05\% & 0.317 \\
No-vine TD3 ensemble & \textbf{62.68\%} & \textbf{30.40\%} & \textbf{12.79\%} & \textbf{2.157} & \textbf{5.213} & 6.75\% & \textbf{29.43\%} & 0.373 \\
Pretrained-only TD3 ensemble & 59.61\% & 29.05\% & 13.02\% & 2.041 & 4.838 & \textbf{6.26\%} & 28.05\% & \textbf{0.314} \\
\bottomrule
\end{tabular}%
}
\caption*{\footnotesize Notes: Post-holdout explanatory evidence on the consumed market path. Bold denotes the best observed value within this three-row comparison and is not a significance claim.}
\end{table}

The full-minus-no-vine mean monthly net-return and utility differences are -0.000824 and -0.000877; the full policy is better in only 8 of 22 periods. The no-vine control also has higher observed CRRA CE and Sharpe. Hence this coarse ablation does not support an incremental benefit from the combined vine-state and scenario-CVaR channels. It does not identify which channel is responsible, nor does non-significance establish equivalence.

The full-minus-pretrained-only mean net-return difference is $3.66\times10^{-5}$ per month and the utility difference is $-2.96\times10^{-7}$. Annualised CRRA CE rounds to 28.05\% for both policies. Historical fine-tuning therefore has no material descriptive increment on this holdout. This finding motivates, rather than substitutes for, the matched-seed component design.

\subsection{Prospective Causal Mechanism Analysis}
\label{subsec:4_causal_results}

The causal analysis contract was frozen before its evaluation output was available. It requires all eight component contrasts and four algorithm controls to be reported regardless of sign or significance, together with raw and Holm-adjusted p-values, annualised certainty-equivalent effects, seed stability, implementation diagnostics, and ensemble wealth paths. It treats the ten matched seeds as optimisation replicates and performs uncertainty estimation over paired market periods.

\IfFileExists{generated/causal_primary_contrasts.tex}{%
\input{generated/causal_primary_contrasts.tex}
}{%
\begin{quote}
\textit{Causal results are intentionally not inserted in this revision because the frozen matched-seed sweep and common-path scorer have not yet produced a validated release. This paragraph must be replaced only from the generated causal tables; no manual or partial-run values are permitted.}
\end{quote}
}

This interface prevents a partially completed training sweep from becoming a manuscript result. Once frozen, the generated section will report the effects of direct vine state, the scenario-CVaR observation, their joint information, CVaR reward shaping, synthetic pre-training, the neural generator relative to temporal bootstrap, recurrent encoding, historical fine-tuning, and the alternative algorithms. Regardless of outcome, these remain post-holdout explanatory findings.

\subsection{Discussion}
\label{subsec:4_discussion}

The frozen evidence supports a mixed conclusion. The proposed ensemble is a competitive low-volatility controller with the highest observed Sharpe ratio and modest implementation drag. It is economically better than equal weight, shrinkage mean--variance, and rolling vine on terminal wealth, and nearly ties DCC--GARCH and the dynamic NN-vine. Yet it loses to the static vine on wealth and CRRA CE, none of its utility effects is statistically significant, and individual-seed behaviour is unstable.

The explanatory ablations sharpen the interpretation. The no-vine ensemble performs better descriptively, while fine-tuning adds almost no utility. Together with the dynamic NN-vine benchmark, these findings mean that performance alone does not establish the novel contribution of the neural vine or RL components. The causal study is necessary to determine whether value comes from the synthetic generator, direct dependence features, scenario risk, recurrent memory, reward shaping, or simply weight-space ensemble averaging.

The study also demonstrates why scientific and operational validation must be integrated. Leakage checks, immutable calendar and asset contracts, behavioural gates, checkpoint metadata, process-isolated inference, constraint audits, and explicit failed-run preservation prevent several common backtest errors. They do not compensate for a short holdout. A credible superiority claim requires new, non-overlapping evidence: frozen expanding or rolling windows, an external investable total-return panel, or future observations arriving after protocol freeze.

\section{Conclusion}
\label{sec:conclusion}

This paper develops a dynamic portfolio framework in which a fully neural, all-tree Student-$t$ D-vine and an LSTM-TD3 controller have distinct but connected roles. The D-vine generates causally ordered synthetic episodes, provides explicit conditional-dependence features, and supplies the scenarios used for CVaR. The agent selects self-financing long--short portfolios under gross and position limits and maximises dense increments of terminal-wealth CRRA utility after transaction and financing costs. The training protocol uses all synthetic data for pre-training, all permissible realised historical trajectories for one fine-tuning pass, and reserves the final 24 decisions for locked evaluation.

The frozen main result is competitive but mixed. Over 22 complete holding periods, the preregistered 20-seed mean-weight ensemble records the lowest observed volatility and highest observed Sharpe ratio among six benchmarks, while ranking third in terminal wealth and CRRA certainty-equivalent return. Its positive utility effect relative to equal weight is not statistically significant, all benchmark-effect intervals include zero, and individual-seed results are dispersed. Weight averaging removes most seed-specific short and leveraged exposure, so part of the ensemble's advantage is diversification and implementation-cost compression rather than a uniformly stable learned policy.

The post-holdout ablations impose an important limit on the novelty claim. A coarse no-vine TD3 ensemble performs better descriptively, and the pretrained-only ensemble is nearly indistinguishable from the fully fine-tuned policy on the consumed path. These results do not prove that vine information or fine-tuning is harmful, because the coarse intervention removes multiple channels and the sample is short. They do show that incremental component value has not yet been established. The preregistered matched-seed causal analysis therefore separates direct vine features, scenario-CVaR observation, CVaR reward shaping, synthetic pre-training, the neural generator, recurrence, fine-tuning, and the TD3 update rule.

Accordingly, the defensible contribution at this stage is methodological and empirical transparency: a dynamic dependence and control system, a validated synthetic-data protocol, publication-grade leakage and constraint controls, and a frozen mixed result that is not converted into a superiority narrative. Confirmation requires evidence independent of the consumed holdout. The next tests should use non-overlapping frozen walk-forward windows, an external investable total-return panel, and ultimately future observations collected after preregistration. Those designs must retain drift-aware turnover, actual-day financing, fail-closed optimisation, paired time-series inference, and complete source, data, environment, and retry provenance.

\appendix
\section{Reproducibility and Causal-Analysis Contract}
\label{sec:appendix}

\subsection{Evidence Classes}

Three evidence classes are kept separate throughout the project. The frozen main evaluation is the original locked out-of-sample result. The no-vine, pretrained-only, ensemble-mechanism, and component analyses are post-holdout explanatory because they reuse its realised path. Retrospective walk-forward and external-panel studies can provide robustness and external-validity evidence, but only future non-overlapping windows frozen before access are eligible for a new confirmatory claim.

\subsection{Causal Contrast Matrix}

\begin{table}[H]
\centering
\caption{Preregistered Post-Holdout Causal Contrasts}
\label{tab:causal_contract}
\resizebox{\textwidth}{!}{%
\begin{tabular}{p{0.31\textwidth} p{0.28\textwidth} p{0.34\textwidth}}
\toprule
\textbf{Contrast} & \textbf{Intervention} & \textbf{Quantity isolated} \\
\midrule
Direct NN-vine state & Zero 63 vine features; retain scenario-CVaR observation & Direct policy value of edge dependence features \\
Scenario-CVaR observation & Retain vine features; zero CVaR state channel & Policy value of the scenario risk statistic \\
Joint visible dependence & Zero vine features and CVaR observation & Combined policy-visible dependence information \\
CVaR reward shaping & Retain state; zero CVaR reward term & Effect of tail-risk shaping in the objective \\
Synthetic pre-training & Historical training only & Contribution of synthetic pre-training \\
Generator comparison & Moving-block-bootstrap pre-training & Neural vine generator versus temporal bootstrap \\
Recurrent encoder & Capacity-matched feedforward encoder & Contribution of LSTM memory \\
Historical fine-tuning & Evaluate pretrained policy only & Contribution of realised-history adaptation \\
Algorithm controls & DDPG, SAC, PPO, and A2C & TD3-specific update-rule sensitivity \\
\bottomrule
\end{tabular}%
}
\end{table}

All variants use seeds 20261001--20261010. A missing or invalid seed invalidates the matched experiment rather than being replaced. Ensembles are constructed by averaging target weights and then recomputing drifted turnover, transaction costs, financing costs, and returns. The primary component family uses one-sided paired CRRA-utility effects with Holm adjustment across eight contrasts; algorithm comparisons are two-sided and exploratory. Circular moving-block bootstrap inference uses 9,999 replications and block length three.

\subsection{Minimum Replication Package}

A complete release should contain: (i) the raw adjusted-level panel or an immutable licensed-data reference and hash; (ii) the training marginal and neural-vine fits; (iii) synthetic and historical bundle manifests; (iv) all seed checkpoints, behaviour gates, schedules, code hashes, and data hashes; (v) frozen training and evaluation source snapshots; (vi) the main and explanatory archives with SHA-256 sidecars; (vii) failed operational-retry logs and an incident timeline; (viii) R, Python, package, CUDA, and container inventories; and (ix) a contents manifest covering every supplementary artifact. Manuscript tables and figures must be generated from those frozen artifacts rather than transcribed manually.
 

% Bibliography
\bibliographystyle{plainnat} 
\bibliography{Zotero}

\end{document}
